\section{Ordered tensor primary packets: invariant characters and the actual arithmetic action}
\label{CFStage:sec:tp-tensor}

This calculation takes the ordered tensor product of the complete
single-primary marked family, with a common parameter \(u\) and with
every original \(t\)-parameter fixed at zero. It does not pass to a
symmetric quotient. Its source is the full packet attached to a selected
actual zero when such a packet is given; it asserts the existence of no
zero off the critical line. The constant in the original primitive,
the original monomial coordinates, and the full source unit are retained.

\subsection{The original ordered family and its period matrix}

Let \(\rho\) be the selected zero of \(g=2\xi\) of complete order
\(m\geq1\), put \(N=m+1\), and let \(k\geq1\) be an integer.
For \(1\leq i\leq k\), put \(y_i=s_i-\rho\). The precise phase and
marked algebra are
\[
 \begin{split}
 h(s)&=(s-\rho)^m,\qquad c=-\frac{(-\rho)^N}{N},\\
 \Phi_k(s_1,\ldots,s_k)
  &=\sum_{i=1}^k\frac{(s_i-\rho)^N-(-\rho)^N}{N}
    =\sum_{i=1}^k\frac{y_i^N}{N}+kc,\\
 E_k&=\mathbb C[s_1,\ldots,s_k]/(h(s_1),\ldots,h(s_k)).
 \end{split}
 \tag{TP.1}\label{CFStage:eq:tp-family}
\]
Each primitive is zero at its original base point \(s_i=0\).
Index every ordered basis lexicographically by
\(\mathbf b=(b_1,\ldots,b_k)\in\{0,\ldots,m-1\}^k\).
Write \(s^{\mathbf b}=\prod_i s_i^{b_i}\),
\(y^{\mathbf b}=\prod_i y_i^{b_i}\),
\(|\mathbf b|=\sum_i b_i\), and
\(L(\mathbf b)=|\mathbf b|+k\).
The unchanged top differential is
\(ds_1\wedge\cdots\wedge ds_k=dy_1\wedge\cdots\wedge dy_k\).

The one-factor Pascal matrix and its ordered tensor power are
\[
 P_{ab}=\begin{cases}\binom ba\rho^{b-a}&a\leq b,\\0&a>b,
             \end{cases}\qquad
 (P^{-1})_{ab}=\begin{cases}\binom ba(-\rho)^{b-a}&a\leq b,\\0&a>b,
             \end{cases}\qquad
 P_k=P^{\otimes k}.
 \tag{TP.2}\label{CFStage:eq:tp-pascal}
\]
Here all one-factor indices run from zero to \(m-1\).
The binomial theorem gives \(e_s=e_yP_k\), where \(e_s,e_y\) are
the ordered rows of monomials. Substitution in both directions proves
the inverse formula, and \(\det P=1\). Coefficient columns therefore
satisfy \(x_y=P_kx_s\), with no change of the polynomial object.

The ordered exponential complex is the tensor product over
\(\mathbb C[u]\) of the two-term complexes with differential
\[
 L_i=u\partial_{s_i}+h(s_i).
 \tag{TP.3}\label{CFStage:eq:tp-differential}
\]
Equivalently its total differential on polynomial forms is
\(u\,d+d\Phi_k\wedge\), with the usual ordered wedge signs.
The operators \(L_i\) commute for different indices, proving that
the square of this differential is zero. In one variable the leading
degree of \(L_iQ\), for \(Q\ne0\), is \(\deg Q+m\), with the
same leading coefficient. Monic subtraction therefore gives a unique
remainder of degree below \(m\). This decomposes the two-term complex
as a contractible complex mapping isomorphically onto its image,
plus a free \(\mathbb C[u]\)-module of rank \(m\) in degree one.
To see contractibility explicitly, send an image element \(L_iQ\)
back to its unique \(Q\), and send its remainder to zero.
Tensoring a contracting homotopy with the identity contracts every
summand having a contractible factor; the total-complex signs cancel
the differentials in the other factors. Consequently all cohomology
vanishes except in degree \(k\), where
\[
 \mathcal H_k=\bigoplus_{\mathbf b}\mathbb C[u]
          [s^{\mathbf b}ds_1\wedge\cdots\wedge ds_k],
 \qquad \mathcal H_k/u\mathcal H_k=E_k\,ds_1\wedge\cdots\wedge ds_k.
 \tag{TP.4}\label{CFStage:eq:tp-lattice}
\]
In particular its rank is \(m^k\), and every primary order is retained.

Fix a branch of \(u^{1/N}\) on a sector of \(u\ne0\), let
\(\zeta=e^{2\pi i/N}\), and for \(0\leq a<m\) define
\[
 \beta_a=\frac{a+1}{N},\quad
 d_a=N^{\beta_a-1}e^{\pi i\beta_a}\Gamma(\beta_a),\quad
 W_{ja}=\zeta^{j(a+1)}-1\ (1\leq j\leq m),\quad
 D(u)=\operatorname{diag}(d_a u^{\beta_a}).
 \tag{TP.5}\label{CFStage:eq:tp-constants}
\]
Use the original rays at angles
\((\arg u+\pi)/N+2\pi j/N\), oriented outwards, and the original
relative contour \(\Gamma_j=-\ell_0+\ell_j\).
The ordered product contour is
\(\Gamma_{j_1}\times\cdots\times\Gamma_{j_k}\), with its product
orientation. There is no permutation of its factors or top form.
The full original period matrix is
\[
 \Pi_k(u)=e^{kc/u}W_kD_k(u)P_k,\qquad
 W_k=W^{\otimes k},\quad D_k(u)=D(u)^{\otimes k}.
 \tag{TP.6}\label{CFStage:eq:tp-period}
\]
Here is a direct verification retaining the original contours.
In each variable, translate its contour to \(\rho+\Gamma_j\).
Cauchy's theorem on the finite quadrilateral swept by each ray
applies to the entire integrand. The junction connectors cancel
between the negative and positive rays. At their remote endpoints
\(s=Re^{i\vartheta}+a\rho\), \(0\leq a\leq1\), the exponent has
real part \(-R^N/(N|u|)+O(R^{N-1})+O(1)\), uniformly in \(a\).
Each connector has length at most \(|\rho|\); its polynomially
weighted integral tends to zero. Absolute integrability on each
ray follows from the same estimate. Translation is thus a proved
contour equality, preserving the scalar \(e^{c/u}\).
On the translated ray the substitution
\(v=r^N/(N|u|)\) gives precisely
\(d_a u^{\beta_a}\zeta^{j(a+1)}\); subtracting the zeroth ray gives
the entry of \(WD(u)\). The Pascal expansion gives the factor \(P\).
Absolute convergence permits Fubini on every product contour and
proves (TP.6), including its orientation and all constants.

The matrix \(W\) is invertible: a relation among its columns gives
a polynomial \(\sum_{r=1}^{m}a_r(z^r-1)\) of degree at most
\(m=N-1\), vanishing at every \(N\)-th root of unity, including
one. A nonzero polynomial cannot have more roots than its degree,
so each \(a_r=0\). Every \(d_a\ne0\), since \(0<\beta_a<1\).
Thus (TP.6) is invertible for every \(u\ne0\). For reference its
determinant, without discarding its scalar, is
\[
 \det\Pi_k(u)=K_m^{k m^{k-1}}
       u^{k m^k/2}e^{k m^k c/u},\qquad
 K_m=N^{-m/2}e^{\pi i m/2}
       \left(\prod_{a=0}^{m-1}\Gamma(\beta_a)\right)\det W.
 \tag{TP.7}\label{CFStage:eq:tp-determinant}
\]
Indeed \(\sum_a\beta_a=m/2\), and
\(\det(X\otimes Y)=(\det X)^{\dim Y}(\det Y)^{\dim X}\).
The latter identity follows first on triangular matrices by
multiplying their diagonal entries and then on all complex matrices
by triangularizing them separately. Applying it repeatedly proves
the displayed exponent \(k m^{k-1}\), also when \(m=1\).

\subsection{Full monodromy and the exact character projectors}

Let \(B_s\) have column \(b\) equal to
\(((b+1)s^b-b\rho s^{b-1})/N\) in the original basis, and let
\(B_k=\sum_i I^{\otimes(i-1)}\otimes B_s\otimes I^{\otimes(k-i)}\).
The exact connection on (TP.4) is
\[
 \nabla_u=\partial_u-\frac{kc}{u^2}I+\frac{B_k}{u},
 \qquad P_kB_kP_k^{-1}
         =\operatorname{diag}(L(\mathbf b)/N).
 \tag{TP.8}\label{CFStage:eq:tp-connection}
\]
For one factor, divide \(\Phi(s)s^b\) monically by \(h(s)\):
the quotient is \((s-\rho)s^b/N\) and the remainder is \(cs^b\).
Subtracting its \(L_u\)-image gives
\([\Phi s^b]=[cs^b-u((s-\rho)s^b/N)']\).
This is exactly \(\partial_u-\Phi/u^2\) on cohomology and gives
the stated column. The binomial coordinate map makes this column
diagonal with entries \(\beta_a\). Summing the tensor-factor
connections proves (TP.8).

On one positive circuit of \(u=0\), each relative contour changes
by \(\Gamma_j\mapsto\Gamma_{j+1}-\Gamma_1\) for \(j<m\), and
\(\Gamma_m\mapsto-\Gamma_1\). This follows directly by advancing
each ray index by one modulo \(N\). Applying this to the columns
of \(W\) gives eigenvalue \(\zeta^{a+1}\).
Consequently the full period-cycle monodromy is
\[
 \Pi_k^{\rm cont}=M_k\Pi_k,\qquad
 M_k=W_k\operatorname{diag}(\zeta^{L(\mathbf b)})W_k^{-1},
 \qquad M_k^N=I.
 \tag{TP.9}\label{CFStage:eq:tp-monodromy}
\]
The scalar \(e^{kc/u}\) is single-valued and has not been removed.
The horizontal coefficient fundamental matrix is
\(P_k^{-1}e^{-kc/u}\operatorname{diag}(u^{-L(\mathbf b)/N})\),
as follows by differentiation using (TP.8). Its monodromy has the
inverse eigenvalues. Both conventions have the same invariant
indices, namely \(N\mid L(\mathbf b)\).

For a residue \(r\in\mathbb Z/N\mathbb Z\), define the diagonal
mask \(E_r\) on the \(y\)-basis by
\(E_r y^{\mathbf b}=\mathbf1_{L(\mathbf b)\equiv r}y^{\mathbf b}\).
The period-space and marked-space projectors are, respectively,
\[
 \begin{split}
 \mathcal P_r&=\frac1N\sum_{j=0}^{N-1}\zeta^{-rj}M_k^j
             =W_kE_rW_k^{-1},\\
 Q_r&=P_k^{-1}E_rP_k,\qquad
       \mathcal P_r\Pi_k=\Pi_k Q_r .
 \end{split}
 \tag{TP.10}\label{CFStage:eq:tp-projector}
\]
The geometric sum of \(N\)-th roots is \(N\) when its exponent
is divisible by \(N\) and zero otherwise; this proves the first
equality and the identities
\(Q_rQ_t=\mathbf1_{r=t}Q_r\), \(\sum_rQ_r=I\), and the same
identities for \(\mathcal P_r\). The diagonal matrix \(D_k(u)\)
commutes with every \(E_r\), proving the last equality. In
particular the pulled-back projector is constant in \(u\).
These are algebraic projectors; no orthogonality in the original
theta-source metric is assumed.

Their dimensions are exactly
\[
 d_r=\operatorname{rank}Q_r
  =\begin{cases}
  \bigl(m^k+m(-1)^k\bigr)/N,&r=0,\\
  \bigl(m^k-(-1)^k\bigr)/N,&r\ne0.
  \end{cases}
 \tag{TP.11}\label{CFStage:eq:tp-count}
\]
In fact the character sum giving the number of ordered tuples is
\(N^{-1}\sum_{j=0}^{N-1}\zeta^{-rj}
(\sum_{l=1}^{N-1}\zeta^{jl})^k\). Its inner sum is \(m\) at
\(j=0\) and \(-1\) otherwise. The sum of the remaining outer
characters is \(m\) for \(r=0\) and \(-1\) for \(r\ne0\).
This proves (TP.11) as an actual integer count, including all
small values of \(m,k\).

The ordinary invariant stalk of this punctured-disc local system
is \(\ker(M_k-I)=\operatorname{im}\mathcal P_0\), of dimension
\(d_0\). Its coinvariant space is identified with that same
summand by \(\mathcal P_0\): on every other summand \(M_k-I\)
is invertible. Thus both circle cohomology groups have dimension
\(d_0\). On the cover \(u=v^N\), (TP.8) becomes
\[
 \partial_v-\frac{Nkc}{v^{N+1}}I
       +\frac1v\operatorname{diag}(L(\mathbf b)).
 \tag{TP.12}\label{CFStage:eq:tp-cover}
\]
The basis change by \(\operatorname{diag}(v^{-L(\mathbf b)})\)
removes precisely the integer residue. The remaining exact
horizontal scalar is \(e^{-kc/v^N}\). It is not a meromorphic
gauge removing the irregular term when \(kc\ne0\).
Nevertheless all displayed solutions are single-valued on this
cover, so its analytic monodromy is the identity and its
unipotent logarithm is zero. The exact diagonal solution leaves
no additional sector transition to determine for this tensor
of the one-critical-value family. The invariant and coinvariant
statements above refer to the ordinary local system; they do not
impose an unmentioned growth condition at the puncture.

\subsection{The arithmetic action moves the invariant character}

Let \(J\) be the matrix of multiplication by \(y\) on
\(\mathbb C[y]/y^m\): \(J_{a+1,a}=1\), all other entries zero.
Put \(J_i=I^{\otimes(i-1)}\otimes J\otimes I^{\otimes(k-i)}\),
\(J_{\rm tot}=\sum_iJ_i\), and
\[
 R_k=P_k^{-1}J_{\rm tot}P_k,\qquad
 A_k=M_{s_1+\cdots+s_k}=k\rho I+R_k.
 \tag{TP.13}\label{CFStage:eq:tp-arithmetic}
\]
This is multiplication in the original algebra (TP.1), since
\(s_i=\rho+y_i\). No geometric nilpotent replaces \(R_k\).
Put \(K=k(m-1)\). The commuting nilpotents have the exact
multinomial powers
\[
 J_{\rm tot}^{r}y^{\mathbf b}
  =\sum_{\substack{\boldsymbol\alpha\in\mathbb Z_{\geq0}^k,
                       |\boldsymbol\alpha|=r\\
                    b_i+\alpha_i\leq m-1\ (1\leq i\leq k)}}
       \frac{r!}{\prod_i\alpha_i!}\,y^{\mathbf b+\boldsymbol\alpha},
 \quad R_k^{K+1}=0,\quad R_k^K\ne0.
 \tag{TP.14}\label{CFStage:eq:tp-powers}
\]
Expansion of the commuting operators proves the sum. A term with
an exponent at least \(m\) vanishes in (TP.1), giving its exact
index restrictions. At \(r=K\), the input one has nonzero top
coefficient \(K!/((m-1)!)^k\); at larger \(r\) no tuple survives.
This also proves the conventions when \(K=0\): \(R_k^0=I\ne0\)
and \(R_k=0\).

Define \(D_0=\operatorname{diag}(d_a)\) and
\(K_0=WD_0JD_0^{-1}W^{-1}\), and let
\(K_{\rm tot}\) be the sum of its ordered factor placements.
The actual period transport of the arithmetic action is
\[
 \Pi_k A_k\Pi_k^{-1}=k\rho I+u^{1/N}K_{\rm tot},
 \qquad M_kK_{\rm tot}M_k^{-1}=\zeta K_{\rm tot}.
 \tag{TP.15}\label{CFStage:eq:tp-transport}
\]
Conjugating \(J\) by \(\operatorname{diag}(u^{\beta_a})\)
multiplies each nonzero entry by
\(u^{\beta_{a+1}-\beta_a}=u^{1/N}\); the tensor placements
give the first identity. Replacing these powers by
\(\zeta^{a+1}\) gives the second. The scalar \(e^{kc/u}\)
cancels by its displayed inverse in this conjugation.
Equivalently the exact connection commutator is
\[
 [B_k,A_k]=R_k/N,\qquad [\nabla_u,A_k]=R_k/(Nu).
 \tag{TP.16}\label{CFStage:eq:tp-commutator}
\]
These follow on each monomial from the difference of its two
residue eigenvalues, including each truncated top term, and then
from (TP.8). Thus the arithmetic action is explicitly related
to the boundary connection, with its precise nonhorizontal term.

Multiplication by a surviving \(y_i\) increases \(L\) by one.
For all residues and integers \(r\geq0\), this gives
\[
 \begin{gathered}
 Q_t R_k^r=R_k^r Q_{t-r},\qquad
 [Q_0,A_k]=R_k(Q_{-1}-Q_0),\\
 Q_0A_kQ_0=k\rho Q_0,\qquad
 (I-Q_0)A_kQ_0=R_kQ_0=Q_1R_kQ_0.
 \end{gathered}
 \tag{TP.17}\label{CFStage:eq:tp-leakage}
\]
All subscripts are taken modulo \(N\). The first identity
holds on every basis monomial, including zero output terms.
The other identities follow from it and from \(N\geq2\).
They calculate both the compression and the outgoing map.
In particular a scalar compression is not an assertion that
the invariant plane is preserved.

For a precise nonvanishing statement put
\[
 d_{\min}=N\lceil k/N\rceil-k\in\{0,\ldots,N-1\}.
 \tag{TP.18}\label{CFStage:eq:tp-minimum}
\]
The possible degrees \(|\mathbf b|\) fill every integer from
zero to \(K\): distribute any such integer successively into
the \(k\) slots, never exceeding their capacity \(m-1\).
Thus \(Q_0\ne0\) exactly when \(d_{\min}\leq K\).
Moreover \(R_kQ_0\ne0\) exactly when \(d_{\min}<K\).
If strict inequality holds, a tuple of that degree has an
unfilled slot, and (TP.14) with \(r=1\) gives nonzero output
with no cancellation, since its distinct monomials have
positive integer coefficients. If equality holds, the only
possible invariant tuple is the top one and every \(J_i\)
kills it. This proves both directions.

For \(k=1\) we have \(d_{\min}=m>K=m-1\), so there are no
invariants. For \(m=1\), \(R_k=0\) and the unique tensor
character is invariant exactly when \(k\) is even.
For every \(m\geq2,k\geq2\), instead,
\[
 Q_0\ne0,\qquad R_kQ_0\ne0.
 \tag{TP.19}\label{CFStage:eq:tp-nonzero}
\]
Indeed at \(k=2\), \(d_{\min}=m-1<2m-2=K\).
For \(k\geq3\), \(K=k(m-1)\geq m+1=N\), while
\(d_{\min}\leq N-1\). This proves (TP.19), including
\(m=2,k=3\), the smallest equality case for the capacity bound.

\subsection{The complete compressed dilation, including every return power}

Work on \(V_0=\operatorname{im}Q_0\). If this space is zero,
all the following compressed maps are the unique zero maps.
Suppose it is nonzero and put
\[
 q_* =\left\lfloor\frac{K-d_{\min}}{N}\right\rfloor
       =\lfloor km/N\rfloor-\lceil k/N\rceil,\qquad
 T=\left.Q_0R_k^NQ_0\right|_{V_0}.
 \tag{TP.20}\label{CFStage:eq:tp-return}
\]
Since \(R_k^N\) commutes with \(Q_0\),
\(T^q=\left.Q_0R_k^{qN}Q_0\right|_{V_0}\).
For every \(0\leq q\leq q_*\) this map is nonzero; for
\(q>q_*\) it is zero. To prove nonvanishing, choose an
invariant tuple of minimum degree. Its remaining capacities
sum to \(K-d_{\min}\), so an increment tuple of size \(qN\)
can be distributed within them. Its term in (TP.14) is
nonzero and cannot cancel. For vanishing, every input has
degree at least \(d_{\min}\), so all output terms exceed
capacity when \(qN>K-d_{\min}\). Thus the nilpotency
index of \(T\) is exactly \(q_*+1\), with the identity
zeroth power when \(q_*=0\).

For \(a>0\), define the original ordered dilation on
functions by simultaneous replacement \(x_i\mapsto x_i/a\).
Its generator is \(\sum_iD_i\), with
\(D_i=-x_i\partial_{x_i}\), and hence its action on this
marked packet is \(a^{A_k}\). Write \(z=\log a\).
The complete compression is the finite expression
\[
 \left.Q_0a^{A_k}Q_0\right|_{V_0}
   =a^{k\rho}\sum_{q=0}^{q_*}
          \frac{z^{qN}}{(qN)!}T^q.
 \tag{TP.21}\label{CFStage:eq:tp-dilation}
\]
Indeed \(a^{A_k}=a^{k\rho}\sum_{r=0}^{K}z^rR_k^r/r!\)
by (TP.14). The first identity in (TP.17) kills its
compression unless \(N\mid r\), and (TP.20) determines
exactly which surviving powers are nonzero. Equivalently,
the sum without its scalar is the restriction of
\(N^{-1}\sum_{j=0}^{N-1}\exp(\zeta^jzR_k)Q_0\);
the same root-of-unity sum proves this equality.
In the original \(y\)-monomial coordinates its exact
entry, from invariant \(\mathbf b\) to invariant
\(\mathbf d\), is
\[
 \begin{cases}
 a^{k\rho}\displaystyle\frac{z^{|\mathbf d|-|\mathbf b|}}
                    {\prod_i(d_i-b_i)!},&d_i\geq b_i\ \hbox{for all }i,\\
 0,&\hbox{otherwise}.
 \end{cases}
 \tag{TP.22}\label{CFStage:eq:tp-entry}
\]
The two invariant conditions make the exponent divisible
by \(N\). Formula (TP.14) cancels its total factorial
against that of the exponential and proves (TP.22).
Conjugation by the explicit columns of \(P_k^{-1}\)
returns this same matrix to the original \(s\)-basis.

In particular \(k=2\) always has \(q_*=0\), so (TP.21)
is scalar even though (TP.19) proves nonzero outgoing
arithmetic action when \(m\geq2\). Conversely nonconstant
returns do occur: for \(m=2,N=3,k=3\), \(V_0\) is spanned
in the \(y\)-basis by \(1\) and \(y_1y_2y_3\), and
\(R_k^3(1)=6y_1y_2y_3\). The compression is exactly
\(a^{3\rho}(I+z^3T/6)\), with that nonzero \(T\).
When \(q_*\geq1\), the compressed maps are not a
one-parameter representation: after removing the scalar,
their derivative at \(z=0\) is zero but their coefficient
of \(z^N\) is the nonzero \(T/N!\). A differentiable
matrix-valued function \(F\) with \(F(0)=I\) and
\(F(z+w)=F(z)F(w)\) satisfies \(F'(z)=F(z)F'(0)\)
by differentiating in \(w\); zero derivative at zero
would make it constant. This proves the asserted failure
without identifying compression with a quotient action.

\subsection{The full Taylor unit and its exact original-source projector}

Retain the entire quotient \(v_h=g/h\) and all coefficients
of its marked Taylor class:
\[
 a_j=\frac{v_h^{(j)}(\rho)}{j!}
     =\frac{g^{(m+j)}(\rho)}{(m+j)!}\quad(0\leq j<m),
 \quad a_0\ne0,\quad
 U_s=P^{-1}\left(\sum_{j=0}^{m-1}a_jJ^j\right)P,
 \quad U_k=U_s^{\otimes k}.
 \tag{TP.23}\label{CFStage:eq:tp-unit}
\]
The Taylor expansion at the complete-order zero proves
the derivative equality. The inverse coefficients are
\(b_0=a_0^{-1}\) and
\(b_r=-a_0^{-1}\sum_{j=1}^ra_jb_{r-j}\); multiplying
the finite series modulo \(J^m\) proves the inverse.
Thus (TP.23) is invertible and keeps every coefficient.
It commutes with \(A_k\), since all its factors are
polynomials in the commuting multiplication matrices.

For a multi-index \(\mathbf j\in\{0,\ldots,m-1\}^k\),
write \(a_{\mathbf j}=\prod_i a_{j_i}\) and let
\(R_i=P_k^{-1}J_iP_k\), so \(R_k=\sum_iR_i\).
The full graded unit, its period transport and its
connection defect are
\[
 \begin{split}
 U_k&=\sum_{\mathbf j}a_{\mathbf j}\prod_iR_i^{j_i},\\
 \Pi_kU_k\Pi_k^{-1}
   &=\sum_{\mathbf j}a_{\mathbf j}u^{|\mathbf j|/N}
                              \prod_iK_i^{j_i},\\
 [\nabla_u,U_k]
   &=\frac1{Nu}\sum_{\mathbf j}|\mathbf j|a_{\mathbf j}
                              \prod_iR_i^{j_i}.
 \end{split}
 \tag{TP.24}\label{CFStage:eq:tp-unittransport}
\]
Here \(K_i\) is the corresponding placement of \(K_0\).
The first equality expands the tensor product, the
second uses (TP.15) factor by factor, and the third
follows from the residue difference \(|\mathbf j|/N\)
on each surviving monomial. Zero terms at the truncation
boundary satisfy the identities as well.

The actual source maps can now be specified. On the
original one-variable spaces retain
\(D=-x\partial_x\),
\(\Theta\phi(x)=2\sum_{r\geq1}\phi(rx)\),
\(\phi_*=(4\pi^2x^4-6\pi x^2)e^{-\pi x^2}\), and
the original \(F_h\) with Mellin transform \(v_h\).
Here the Mellin convention is
\(\mathcal MF(s)=\int_0^\infty F(x)x^s\,dx/x\).
The source \(V\) is the even Schwartz subspace with
\(\phi(0)=\int\phi=0\); the target \(\mathscr B\)
has all seminorms \(\sup_{x>0}x^b|D^rF(x)|\) finite
for every integer \(b\) and nonnegative integer \(r\).
The inherited original source identity is
\(\mathcal M\Theta\phi_*=g\), and the exact Euler inverse
constructs this \(F_h\in\mathscr B\).
Define the actual finite-packet maps on multivariate
polynomials by
\[
 \begin{split}
 \mathcal T_k(f)&=f(D_1,\ldots,D_k)
                       \prod_{i=1}^k F_h(x_i),\\
 \alpha_k(f)&=f(D_1,\ldots,D_k)
                       \prod_{i=1}^k\phi_*(x_i),\\
 \mathcal J_k\mathcal T_k(f)&=U_k[f]_s.
 \end{split}
 \tag{TP.25}\label{CFStage:eq:tp-source}
\]
Here \(\mathcal J_k\) means the multivariable Mellin
transform followed by reduction modulo all \(h(s_i)\),
and \([f]_s\) is the unchanged remainder column.
Every displayed source lies in the original smooth
product source; it is an explicit finite sum of tensor
functions, not an assertion that the whole completed
source equals its algebraic tensor product.
Integration by parts gives
\(\mathcal M(D_iF)=s_i\mathcal MF\); all endpoint
terms vanish by the stated seminorms. Fubini then gives
\(\mathcal M_k\mathcal T_k(f)=f(s)\prod_i v_h(s_i)\).
Its full jet is exactly the last equality of (TP.25).
The factorwise identities
\(\Theta\alpha_h(f)=\mathcal T_h(hf)\)
give the ordered tensor source square with every usual
total-complex sign. The corrected one-factor section
\(\sigma_h\) therefore gives the actual degree-\(k\)
injection \(\sigma_h^{\otimes k}U_k\) on this finite
tensor packet. This retains the original full unit;
it is not set equal to the identity.

In particular, if \(x_s\) is a marked coefficient column
and \(q=U_kx_s\) is its actual theta-jet coordinate,
the period observation of that source is
\(\Pi_kU_k^{-1}q\). Its exact invariant projector is
\[
 \widehat Q_0=U_kQ_0U_k^{-1},\qquad
 \mathcal P_0\Pi_kU_k^{-1}
      =\Pi_kU_k^{-1}\widehat Q_0,\qquad
 [\widehat Q_0,A_k]=U_k[Q_0,A_k]U_k^{-1}.
 \tag{TP.26}\label{CFStage:eq:tp-sourceprojector}
\]
The first two formulas follow by substituting
\(q=U_kx_s\) into (TP.10). The last uses the proved
commutation of \(U_k\) with \(A_k\).
Thus the actual invariant source plane is
\(U_k\operatorname{im}Q_0\), with the same exact
arithmetic leakage conjugated through its original unit.
On the cohomology image the corresponding projector is
\(\sigma_h^{\otimes k}\widehat Q_0
(\sigma_h^{\otimes k})^{-1}\), where the inverse is
only on that injective image. Its maps are fully specified.

For comparison, when \(V_0\ne0\), compressing the coefficient matrix
\(U_k\) by \(Q_0\) is another concrete operation. Its
full expression on \(V_0\) is
\[
 \left.Q_0U_kQ_0\right|_{V_0}
  =a_0^k I_{V_0}+Z,\qquad
 Z=\left.\sum_{\substack{\mathbf j\ne0\\N\mid|\mathbf j|}}
       a_{\mathbf j}Q_0\prod_iR_i^{j_i}Q_0\right|_{V_0},
 \qquad Z^{q_*+1}=0.
 \tag{TP.27}\label{CFStage:eq:tp-unitcompression}
\]
Every omitted term changes the character and so is
killed by (TP.10); every term of \(Z\) raises degree
by at least \(N\). The degree bound defining \(q_*\)
proves its stated vanishing power, without requiring
any \(a_j\) other than \(a_0\) to be nonzero.
Its exact inverse is
\(a_0^{-k}\sum_{r=0}^{q_*}(-a_0^{-k}Z)^r\).
Multiplying this finite geometric series proves the
claim. This inverse is not substituted for the full
\(U_k^{-1}\) in (TP.26).

Finally the difference between marked coefficient
multiplication and multiplication of twisted representatives
is also explicit. Put
\(v_{\rm pol}(s)=\sum_{j=0}^{m-1}a_j(s-\rho)^j\)
and \(V_{\rm pol}=\prod_i v_{\rm pol}(s_i)\).
For every polynomial \(f\), the actual differential obeys
\[
 V_{\rm pol}L_if=L_i(V_{\rm pol}f)
                    -u(\partial_{s_i}V_{\rm pol})f.
 \tag{TP.28}\label{CFStage:eq:tp-unitdefect}
\]
This is the Leibniz rule, with the full derivative term.
To compute its complete remainder, divide in one variable
\(v_{\rm pol}s^b=hq_b+r_b\), with \(\deg r_b<m\).
Let \(H_s\) have column \(b\) the original coordinate
column of \(q_b'\). Since \(\deg q_b\leq m-2\),
\(q_b'\) is already a remainder. Subtracting
\(L_uq_b=hq_b+uq_b'\) gives
\([v_{\rm pol}s^b]=[r_b-uq_b']\).
For explicit coefficients in the \(y\)-basis,
\[
 \begin{gathered}
 (H_y)_{ab}=
 \begin{cases}
 (a+1)a_{a+m+1-b},&0\leq a+m+1-b<m,\\
 0,&\hbox{otherwise},
 \end{cases}
 \quad H_s=P^{-1}H_yP,\\
 \operatorname{Rem}_u(V_{\rm pol}\,\cdot)
           =(U_s-uH_s)^{\otimes k}.
 \end{gathered}
 \tag{TP.29}\label{CFStage:eq:tp-unitremainder}
\]
on the retained original remainder frame. Indeed the
quotient of \(v_{\rm pol}y^b\) has terms
\(a_jy^{j+b-m}\) exactly when \(j+b\geq m\);
differentiating gives the stated matrix entry.
The tensor formula follows by replacing the factors
successively; their variables are distinct and their
top-form order stays fixed. This proves the exact full
representative correction on the same frame. It does
not assert that multiplication descends on arbitrary
twisted-cohomology representatives.

\subsection{The original terminal class on its exact invariant subsequence}

The complete tensor primary algebra contains its original terminal
jet
\[
 \tau_k=\prod_{i=1}^k y_i^{m-1},\qquad
 R_k\tau_k=0,\qquad A_k\tau_k=k\rho\tau_k,\qquad
 U_k\tau_k=a_0^k\tau_k.
 \tag{TP.30}\label{CFStage:eq:tp-terminal}
\]
Its displayed coefficient in the retained basis is one, so it is
nonzero. Each multiplication by \(y_i\) kills it, proving the next
two statements by (TP.13). Every positive Taylor term in (TP.23)
also kills it, proving the last statement with its full scalar.
Since \(L(m-1,\ldots,m-1)=km\equiv-k\pmod N\),
\[
 Q_0\tau_k=\widehat Q_0\tau_k
   =\begin{cases}\tau_k,&N\mid k,\\0,&N\nmid k.\end{cases}
 \tag{TP.31}\label{CFStage:eq:tp-terminalprojector}
\]
The equality for \(Q_0\) is its diagonal-mask definition, and
that for \(\widehat Q_0\) follows from
\(U_k^{-1}\tau_k=a_0^{-k}\tau_k\) in (TP.26).
Thus on the infinite subsequence \(k=N\ell\), \(\ell\geq1\),
the same nonzero arithmetic terminal class lies in the actual
boundary-invariant plane, in both marked and original jet coordinates.

Its complete period vector is
\[
 \Pi_k\tau_k=e^{kc/u}d_{m-1}^k u^{km/N}
       (W_{\bullet,m-1})^{\otimes k},\qquad
 d_{m-1}=N^{-1/N}e^{\pi i m/N}\Gamma(m/N),\quad
 W_{j,m-1}=\zeta^{-j}-1.
 \tag{TP.32}\label{CFStage:eq:tp-terminalperiod}
\]
This is column \((m-1,\ldots,m-1)\) of (TP.6) in its
specified \(y\)-coordinates. Every entry \(W_{j,m-1}\)
is nonzero, so the vector remains nonzero for \(u\ne0\).
Its period-cycle monodromy is multiplication by
\(\zeta^{km}=\zeta^{-k}\), agreeing with (TP.31).
For the original corrected-section class
\(u_\rho^{\otimes k}=\sigma_h^{\otimes k}\tau_k\), the
period observation defined in (TP.26) is precisely
\(a_0^{-k}\Pi_k\tau_k\), not the vector with this scalar omitted.

For the comparison with the original MAI/PAM counterfactual branch,
retain its selected \(\rho=1/2+\delta+i\gamma\), with
\(0<\delta<1/2\) and \(\gamma>0\), reached by its original
reflection and conjugation arrows. Let \(h_0\) be its original
reflected quartet packet containing the same order-\(m\) zero, put
\(w=h_0/h\), and retain its original CRT idempotent \(e_\rho\).
The map \(E_h\to e_\rho E_{h_0}\), \([f]\mapsto e_\rho[f]\),
is an isomorphism with inverse reduction modulo \(h\):
the idempotent is one on this primary factor and zero on all
the other factors. This proves both composites directly on
every CRT component. Moreover \(w(\rho)\ne0\), and the exact
unit relation on this factor is
\((g/h_0)=v_h/w\); the inverse of the class of \(w\) is its
unique truncated Taylor inverse. Thus this local passage retains
the whole original unit, not only its constant coefficient.

Write \(K=k(m-1)\) as before. The original amplified terminal
coefficient is
\[
 c_{h_0,k}=\frac{K!}{((m-1)!)^k}
             \bigl((g/h_0)(\rho)\bigr)^k\ne0,
 \qquad \eta_k v_k=c_{h_0,k}\tau_k,
 \qquad \mathcal I_kv_k=c_{h_0,k}u_\rho^{\otimes k}.
 \tag{TP.33}\label{CFStage:eq:tp-MAIscalar}
\]
Here \(v_k,\eta_k,\mathcal I_k\) are the original cyclic
terminal vector, coefficient map and cohomology map of MAI6a.
The selected quartet corner \(k\rho\) arises only by choosing
\(\rho\) in every tensor factor, as in that construction.
The first equality for that map is also its direct local
calculation: (TP.14) at \(r=K\) leaves precisely
\(K!/((m-1)!)^k\prod_i y_i^{m-1}\); the original complete unit
acts on that terminal monomial by \(((g/h_0)(\rho))^k\).
All other Taylor terms vanish because they exceed a retained
primary order. The original corrected-section transition
\(\sigma_{h_0}\iota_{h,h_0}=\sigma_h\) carries this identity
to the second cohomology equality. Thus the original source
period observation of this specified amplified vector is
exactly
\[
 c_{h_0,k}a_0^{-k}\Pi_k\tau_k
 =\frac{K!}{((m-1)!)^k}
       \left(\frac{(g/h_0)(\rho)}{(g/h)(\rho)}\right)^k
       e^{kc/u}d_{m-1}^k u^{km/N}(W_{\bullet,m-1})^{\otimes k}.
 \tag{TP.34}\label{CFStage:eq:tp-MAIperiod}
\]
Every factor is retained and nonzero over \(\mathbb C\).
The ratio of unit values is \(w(\rho)^{-1}\), as follows
from the full unit relation above. No rescaling of the
original vector was used to get (TP.34).

For \(a>0\), (TP.30) gives
\(a^{A_k}\tau_k=a^{k\rho}\tau_k\); the original source
action has the same equality after applying its established
theta-boundary homotopy and passing to cohomology. If
\(\delta=\operatorname{Re}\rho-1/2\), its exact discrepancy is
\[
 \log\!\left(\frac{|a^{k\rho}|^2}{a^k}\right)
        =2k\delta\log a.
 \tag{TP.35}\label{CFStage:eq:tp-terminaldiscrepancy}
\]
This follows by taking the real part of \(k\rho\log a\).
On \(k=N\ell\), it holds on the same nonzero original
terminal class inside the boundary-invariant plane.
The invariant projection therefore does not eliminate
that discrepancy. This calculation supplies no bound
on it by a constant independent of \(k\).

\subsection{The exact rank of arithmetic departure from the invariant plane}

Let \(E_{k,d}\) be the span of \(y^{\mathbf b}\) with
\(|\mathbf b|=d\), and define \(h_d=\dim E_{k,d}\), with
\(h_d=0\) outside \(0\leq d\leq K\). The exact coefficient
count is
\[
 \sum_{d=0}^K h_d t^d=(1+t+\cdots+t^{m-1})^k.
 \tag{TP.36}\label{CFStage:eq:tp-hilbert}
\]
Expanding the product proves it by the ordered monomial
basis. To compute the rank of \(J_{\rm tot}\) on these
levels, define on one factor
\(Fy^b=b(m-b)y^{b-1}\), with \(F1=0\), and use the
positive diagonal form with
\(\|y^b\|^2=b!(m-1)!/(m-1-b)!\).
The ratio of consecutive diagonal entries is
\((b+1)(m-1-b)\); hence the adjoint of \(J\) is exactly
\(F\). This form is used only to prove a basis-independent
algebraic rank. It is not substituted for the original
theta-source norm.

Use the tensor product of these forms and put
\(F_{\rm tot}=\sum_iF_i\). Direct evaluation on each
one-factor monomial gives \([F,J]y^b=(m-1-2b)y^b\).
Operators in different factors commute, so on degree
\(d\) this proves
\[
 \|J_{\rm tot}x\|^2-\|F_{\rm tot}x\|^2
       =(K-2d)\|x\|^2\qquad(x\in E_{k,d}).
 \tag{TP.37}\label{CFStage:eq:tp-rankidentity}
\]
If \(d<K/2\), the identity makes \(J_{\rm tot}\)
injective on \(E_{k,d}\). If \(d+1>K/2\), its reversed
identity makes \(F_{\rm tot}\) injective on
\(E_{k,d+1}\), so its adjoint
\(J_{\rm tot}:E_{k,d}\to E_{k,d+1}\) is surjective:
\((\operatorname{im}J_{\rm tot})^\perp
  =\ker(F_{\rm tot}|_{E_{k,d+1}})=0\).
Every integer \(d\) satisfies
one or both of these inequalities, including the
central degree. It follows that
\[
 \operatorname{rank}(J_{\rm tot}:E_{k,d}\to E_{k,d+1})
         =\min(h_d,h_{d+1}).
 \tag{TP.38}\label{CFStage:eq:tp-rank}
\]
No unimodality assumption is used; the necessary
dimension inequalities follow from the proved
injectivity or surjectivity.

The invariant degrees have \(d+k\equiv0\pmod N\).
Their outgoing images have distinct degrees, so the
direct-sum decomposition and (TP.38) give the full
departure rank and its exact kernel:
\[
 \begin{split}
 \operatorname{rank}(R_kQ_0)
    &=\sum_{\substack{0\leq d\leq K\\N\mid d+k}}
                      \min(h_d,h_{d+1}),\\
 \dim(\ker R_k\cap\operatorname{im}Q_0)
    &=\sum_{\substack{0\leq d\leq K\\N\mid d+k}}
                      \max(h_d-h_{d+1},0).
 \end{split}
 \tag{TP.39}\label{CFStage:eq:tp-leakagerank}
\]
Conjugation by \(P_k\), then by the original \(U_k\)
in (TP.26), preserves these ranks and kernels on the
actual marked and original-source invariant planes.

For example, when \(k=2,m\geq2\), the sole invariant
degree is \(m-1\), and the exact kernel is the line
\[
 \mathbb C\,\omega_m,\qquad
 \omega_m=\sum_{i=0}^{m-1}(-1)^i y_1^i y_2^{m-1-i},
 \qquad (y_1+y_2)\omega_m=0\quad\hbox{in }E_2.
 \tag{TP.40}\label{CFStage:eq:tp-twofactor}
\]
The interior coefficients cancel pairwise, and the two
end terms have exponent \(m\), so vanish. For a general
input \(\sum_i c_i y_1^i y_2^{m-1-i}\), the output
coefficients are \(c_{j-1}+c_j\), \(1\leq j<m\).
Its kernel therefore has exactly the stated alternating
coefficients and dimension one, and its rank is \(m-1\).
This is a nonzero arithmetic eigenclass of eigenvalue
\(2\rho\) already lying in the boundary-invariant plane.
It coexists with the nonzero departure on the remaining
\(m-1\) directions and with the scalar compressed
dilation of (TP.21).

Every constant coefficient map above acts at each original
nonbottom split-support label by \((A,x)\mapsto(A,fx)\),
and fixes the external absence symbol \(\tau\). Its killed
coefficient remains the represented zero at that label.
The marked specialization \(u=0\) is the original
represented-zero specialization; none of these projector
computations replaces it by external absence. Equations
(TP.10), (TP.17), (TP.21) and (TP.26) specify the exact maps
relating boundary invariants, the full arithmetic action,
its dilation and the original source. The dimension of
that invariant space alone supplies no conclusion about
the location of the selected zero.

% BEGIN COMPLETE TENSOR COMPANIONS
% BEGIN TPF SHA256 5c2fa0e144f74aea4f21604e986a36fd8b5f37ee4e8235a36a605b49eba32b95
\section{Ordered tensors of the actual single-primary finite-field sheaf}
\label{CFStage:sec:TPF}

\subsection{The original coefficient map and the common parameter}

Retain the specified coefficient map \(R\to k_0=\mathbf F_{Q_0}\)
of characteristic \(p\), the full Taylor class
\(\upsilon_h=j_h(g/h)\) and its inverse in the retained
\(m\)-dimensional remainder basis, and every other original
coefficient. Write
\[
 N=m+1\ge2,\quad p>N,\qquad
 h(s)=(s-\rho)^m,\qquad
 \Phi(s)=\frac{(s-\rho)^N-(-\rho)^N}{N},\qquad
 c=-\frac{(-\rho)^N}{N}.
 \tag{TPF.1}\label{CFStage:eq:TPF1}
\]
The same coefficient map supplies \(\rho,c,N^{-1}\). In
particular, no coefficient of an infinite analytic function
is being adjoined in place of its specified full Taylor class.
Set
\[
 d=\operatorname{ord}_{(\mathbf Z/N\mathbf Z)^\times}(Q_0),
 \qquad L=\mathbf F_Q,\quad Q=Q_0^d,\quad N\mid Q-1.
 \tag{TPF.2}\label{CFStage:eq:TPF2}
\]
This is defined since \(p>N\) implies \(\gcd(Q_0,N)=1\).
Fix \(\ell\ne p\), coefficients \(\overline{\mathbf Q}_\ell\),
and a nontrivial additive character \(\psi_0:k_0\to
\overline{\mathbf Q}_\ell^\times\). Put
\(\psi=\psi_0\circ\operatorname{Tr}_{L/k_0}\).
Every Frobenius below is geometric Frobenius. Kummer sheaves
\(\mathcal K_\chi\) have trace \(\chi\) on \(L^\times\);
over finite extensions their characters extend by norm.
Artin--Schreier characters extend by trace.

Let \(\pi:\mathbf A^1_s\times\mathbf G_{m,u}\to\mathbf
G_{m,u}\) be projection and retain
\[
 \mathcal F_0=R^1\pi_!\mathcal L_{\psi_0}(\Phi(s)/u),\qquad
 \mathcal T_{0,k}=\mathcal F_0^{\otimes k},\qquad k\ge1.
 \tag{TPF.3}\label{CFStage:eq:TPF3}
\]
The positive integer \(k\) has no restriction relative to
\(p\) or \(N\). The tensor is ordered and taken on this
one common \(u\)-line. No symmetric, alternating, or
sum-coordinate quotient is taken. Write
\(\mathcal F=(\mathcal F_0)_L\) and
\(\mathcal T_k=(\mathcal T_{0,k})_L\).
The complete single-primary calculation SPF gives
\[
 \mathcal F\simeq
 \mathcal L_\psi(c/u)\otimes
 \bigoplus_{\substack{\chi\in X_N\\\chi\ne1}}
        \mathcal A_{a_\chi}\otimes\mathcal K_\chi(u),
 \quad
 X_N=\{\chi:L^\times\to\overline{\mathbf Q}_\ell^\times:
                      \chi^N=1\},
 \tag{TPF.4}\label{CFStage:eq:TPF4}
\]
where all constants are exactly
\[
 G_L(\chi,\psi)=\sum_{z\in L^\times}\chi(z)\psi(z),
 \qquad
 a_\chi=-\chi(N)G_L(\chi,\psi).
 \tag{TPF.5}\label{CFStage:eq:TPF5}
\]
Here \(\mathcal A_a\) is the geometrically constant line
with geometric \(\operatorname{Fr}_Q\)-eigenvalue \(a\).
SPF proves this isomorphism using the finite power-map
idempotents, the support-preserving translation
\(s=y+\rho\), and \(w=z/(Nu)\). Its minus sign comes from
compact cohomology in degree one. In particular, all
\(R^i\pi_!\) vanish for \(i\ne1\). We use that proved
single-primary result without changing its source or phase.

The relative compact-support Kunneth map consequently gives
the precise source realization
\[
 \mathcal T_{0,k}\simeq
 R^k\Pi_!\mathcal L_{\psi_0}
       \left(\frac{\sum_{i=1}^k\Phi(s_i)}u\right),
 \quad
 \Pi:\mathbf A^k_{s_1,\ldots,s_k}\times\mathbf G_{m,u}
          \longrightarrow\mathbf G_{m,u}.
 \tag{TPF.6}\label{CFStage:eq:TPF6}
\]
Indeed the derived relative Kunneth isomorphism tensors the
\(k\) complexes \(R\pi_!\mathcal L_{\psi_0}(\Phi/u)\),
each concentrated in degree one; their tensor is
concentrated in degree \(k\). The external product of the
additive sheaves is the sheaf of the sum of their functions,
by addition of Artin--Schreier torsor coordinates. No
factors are reordered, so no permutation sign is inserted.
The source theorem is SGA~4, Expose XVII, 5.4.3
and its finite-family version 5.4.4.1: apply it
at the finite coefficient levels of these
lisse \(\ell\)-adic sheaves and then pass to
the \(\ell\)-adic system and invert \(\ell\).
Here every projection is compactifiable by
\(\mathbf P^1\times\mathbf G_m\to\mathbf G_m\);
the product compactification is
\((\mathbf P^1)^k\times\mathbf G_m\).
The coordinate isomorphism on this same source is
\[
 (y_1,\ldots,y_k,u)\longmapsto
       (y_1+\rho,\ldots,y_k+\rho,u),\qquad
 \frac{\sum_i\Phi(s_i)}u
       =\frac{\sum_i y_i^N}{Nu}+\frac{kc}{u}.
 \tag{TPF.7}\label{CFStage:eq:TPF7}
\]
Its inverse subtracts \(\rho\) in every coordinate, and it
preserves compact supports. Each original coefficient
function is pulled back by this literal substitution,
and its inverse is pulled back by the same map.

\subsection{Every ordered character tuple and its exact constant}

Define
\[
 \begin{aligned}
 \Omega_k&=(X_N\setminus\{1\})^k,&
 r_k&=(N-1)^k,\\
 \eta_{\boldsymbol\chi}&=\prod_{i=1}^k\chi_i,&
 A_{\boldsymbol\chi}&=\prod_{i=1}^k
            \bigl(-\chi_i(N)G_L(\chi_i,\psi)\bigr),&
 w_k&=kc\in k_0 .
 \end{aligned}
 \tag{TPF.8}\label{CFStage:eq:TPF8}
\]
Thus \(k\) is reduced in characteristic \(p\) only when it
multiplies the field element \(c\). Its integer value still
counts all \(k\) tensor factors and remains in every
cohomological sign. Distributing the ordered tensor in
\eqref{CFStage:eq:TPF4}, and using the additive and multiplicative
torsor tensor laws, proves
\[
 \boxed{\displaystyle
 \mathcal T_k\simeq
 \mathcal L_\psi(w_k/u)\otimes
   \bigoplus_{\boldsymbol\chi\in\Omega_k}
        \mathcal A_{A_{\boldsymbol\chi}}
                  \otimes\mathcal K_{\eta_{\boldsymbol\chi}}(u).}
 \tag{TPF.9}\label{CFStage:eq:TPF9}
\]
Different ordered tuples remain different summands, even
when their product characters or constants agree.
This sheaf is lisse of rank \(r_k\).

For \(L_a=\mathbf F_{Q^a}\), put
\(\psi_a=\psi\circ\operatorname{Tr}_{L_a/L}\) and
\(\chi_{i,a}=\chi_i\circ\operatorname{N}_{L_a/L}\).
At \(u\in L_a^\times\), the eigenvalue on the tuple summand
is exactly
\[
 \psi_a(w_k/u) A_{\boldsymbol\chi}^{\,a}
              \prod_i\chi_{i,a}(u).
 \tag{TPF.10}\label{CFStage:eq:TPF10}
\]
SPF proves \(-G_{L_a}(\chi_{i,a},\psi_a)
=(-G_L(\chi_i,\psi))^a\) by its one-dimensional compact
cohomology. Also \(\chi_{i,a}(N)=\chi_i(N)^a\).
These two equalities prove the stated extension-field
constant, without replacing a negative Gauss sum by
the negative of its \(a\)th power. The trace identity on
the actual ordered source is consequently
\[
 \operatorname{Tr}(\operatorname{Fr}_{Q^a,u}\mid\mathcal T_k)
 =(-1)^k\sum_{(s_1,\ldots,s_k)\in L_a^k}
       \psi_a\left(\frac{\sum_i\Phi(s_i)}u\right).
 \tag{TPF.11}\label{CFStage:eq:TPF11}
\]
To check its sign directly, the original sum factors into
\(k\) one-variable sums, each the negative of the trace on
\(\mathcal F\). Multiplying gives precisely \((-1)^k\).

\subsection{Original inertia, its tame projector, and multiplicities}

Let \(I_0\) be geometric inertia over \(\bar k_0((u))\),
and put \(V_k=(\mathcal T_{0,k})_{\bar\eta}\).
Choose \(\vartheta:\mathbf F_p\to
\overline{\mathbf Q}_\ell^\times\) nontrivial and the unique
\(b\in k_0^\times\) satisfying
\[
 \psi_0(x)=\vartheta(\operatorname{Tr}_{k_0/\mathbf F_p}(bx)).
 \tag{TPF.12}\label{CFStage:eq:TPF12}
\]
Nondegeneracy of the finite-field trace pairing supplies
this unique \(b\); after extension to \(L\) the same \(b\)
represents \(\psi\).

First let \(c\ne0\). Keep the original cover for
\(\mathcal F\), before taking a tensor:
\[
 u=v^N,\qquad z^p-z=\frac{bc}{v^N}.
 \tag{TPF.13}\label{CFStage:eq:TPF13}
\]
The Kummer degree is \(N\). Over \(\bar k_0((v))\) the
Artin--Schreier polynomial has no root, since its right
side has pole order \(N\), prime to \(p\), while a pole
of \(f^p-f\) has order divisible by \(p\). Its degree is
therefore \(p\). The total extension has degree \(Np\),
and the \(Np\) displayed automorphisms
\[
 (\zeta,a):(v,z)\longmapsto(\zeta v,z+a),
 \quad \zeta\in\mu_N,\ a\in\mathbf F_p
 \tag{TPF.14}\label{CFStage:eq:TPF14}
\]
commute and exhaust its Galois group
\(\Gamma=\mu_N\times\mathbf F_p\). This is a quotient of
the original geometric inertia: the Kummer part is
totally ramified, and if \(e\) is the ramification index
of the Artin--Schreier extension over \(\bar k_0((v))\),
its defining equation gives \(p\,v_E(z)=-Ne\).
Since \(\gcd(p,N)=1\) and \(1\le e\le p\), it follows
that \(e=p\).

For \(\chi\in X_N\), let
\(\lambda_\chi:\mu_N(L)\to\overline{\mathbf Q}_\ell^\times\)
be uniquely specified by
\(\lambda_\chi(x^{(Q-1)/N})=\chi(x)\).
Under the geometric Frobenius convention fixed above, the
exact action on the tuple line in \eqref{CFStage:eq:TPF9} is
\[
 \rho_k(\zeta,a)\big|_{\boldsymbol\chi}
    =\lambda_{\eta_{\boldsymbol\chi}}(\zeta)^{-1}
                         \vartheta(-ka).
 \tag{TPF.15}\label{CFStage:eq:TPF15}
\]
The Kummer inverse is the power-map projector convention
of SPF: with \(D_\zeta e_r=e_{\zeta r}\), its
\(\chi\)-projector is
\(N^{-1}\sum_\zeta\lambda_\chi(\zeta)D_\zeta\), whose
image has eigencharacter \(\lambda_\chi^{-1}\).
For the Artin--Schreier torsor \(z^p-z=f\), the
projector \(p^{-1}\sum_a\vartheta(a)D_a\) has
eigencharacter \(\vartheta^{-1}\).
At a finite-field point geometric Frobenius translates
the torsor coordinate by
\(-\operatorname{Tr}(f)\), giving trace
\(\vartheta(\operatorname{Tr}(f))\).
Thus its action is \(\vartheta(-a)\), and \(k\) ordered
tensor factors give \(\vartheta(-ka)\), proving every
sign in \eqref{CFStage:eq:TPF15}.

When \(c=0\), the original additive torsor is trivial.
The geometric inertia action instead factors through
the Kummer quotient \(\mu_N\), with the exact action
\(\lambda_{\eta_{\boldsymbol\chi}}^{-1}\); there is no
nontrivial wild factor.

Let \(T_\zeta\) denote the action of \((\zeta,0)\) in
\eqref{CFStage:eq:TPF15}, or its Kummer action when \(c=0\).
Its exact tame projector is
\[
 P^{\mathrm{tame}}_{0,k}=\frac1N\sum_{\zeta\in\mu_N}T_\zeta,
 \qquad
 P^{\mathrm{tame}}_{0,k}\big|_{\boldsymbol\chi}
   =\begin{cases}
     1,&\eta_{\boldsymbol\chi}=1,\\
     0,&\eta_{\boldsymbol\chi}\ne1 .
    \end{cases}
 \tag{TPF.16}\label{CFStage:eq:TPF16}
\]
Character orthogonality proves the second formula and
idempotence. In the case \(c\ne0\), the true original
inertia projector is
\[
 \begin{split}
 P_{I_0,k}
 &=\frac1{Np}\sum_{\zeta\in\mu_N}\sum_{a\in\mathbf F_p}
                         \rho_k(\zeta,a)\\
 &=P^{\mathrm{tame}}_{0,k}
                \left(\frac1p\sum_{a\in\mathbf F_p}
                                      \vartheta(-ka)\right)
 =\begin{cases}
    P^{\mathrm{tame}}_{0,k},&p\mid k,\\
    0,&p\nmid k .
   \end{cases}
 \end{split}
 \tag{TPF.17}\label{CFStage:eq:TPF17}
\]
Division by \(Np\) occurs in the characteristic-zero
coefficient field, regardless of the prime \(\ell\).
For \(c=0\), the true projector is
\(P_{I_0,k}=P^{\mathrm{tame}}_{0,k}\).
This calculates it inside the original representation;
the tame projector alone is not substituted for full
inertia when its wild character is nontrivial.

Here are all tame multiplicities. Identify the cyclic
group \(X_N\) with the additive group \(\mathbf Z/N\mathbf Z\)
using a chosen generator. A tuple is an ordered sequence
\((j_1,\ldots,j_k)\), each \(1\le j_i<N\).
For \(t\in\mathbf Z/N\mathbf Z\), let \(D_{k,t}\) count
those with \(\sum_i j_i=t\). Choose any primitive complex
\(N\)th root \(\omega\). The elementary root-of-unity
filter yields
\[
 \begin{split}
 D_{k,t}
 &=\frac1N\sum_{a=0}^{N-1}\omega^{-at}
                  \left(\sum_{j=1}^{N-1}\omega^{aj}\right)^k\\
 &=\begin{cases}
 \displaystyle\frac{(N-1)^k+(N-1)(-1)^k}{N},&t=0,\\[4pt]
 \displaystyle\frac{(N-1)^k-(-1)^k}{N},&t\ne0 .
 \end{cases}
 \end{split}
 \tag{TPF.18}\label{CFStage:eq:TPF18}
\]
The filter follows by summing
\(N^{-1}\sum_a\omega^{a(\sum j_i-t)}\), which is one
when the congruence holds and zero otherwise.
At \(a=0\) the inner sum is \(N-1\); for every other
\(a\) it is \(-1\). Finally
\(\sum_{a=1}^{N-1}\omega^{-at}\) is \(N-1\) for \(t=0\)
and \(-1\) otherwise. This proves the count, including
all repeated tuple occurrences.

Put \(d_{0,k}=D_{k,0}\). Since a field has no nonzero
nilpotents and \(N\ne0\), \(c=0\) holds exactly when
\(\rho=0\) after the given specialization. Consequently
\[
 w_k=0\ \Longleftrightarrow\ (c=0\ \hbox{or}\ p\mid k),
 \qquad
 \boxed{\displaystyle
 \dim V_k^{I_0}=
 \begin{cases}
 d_{0,k}=\dfrac{(N-1)^k+(N-1)(-1)^k}{N},&w_k=0,\\[5pt]
 0,&w_k\ne0 .
 \end{cases}}
 \tag{TPF.19}\label{CFStage:eq:TPF19}
\]
This allows \(d_{0,k}=0\): for \(k=1\) it is zero,
and for \(N=2\) it is zero at odd \(k\), one at even
\(k\). For \(N\ge3\) and \(k\ge2\), both formulas
in \eqref{CFStage:eq:TPF18} are strictly positive: at even
\(k\), \((N-1)^k>1\); at odd \(k\ge3\),
\((N-1)^k>N-1\).

The action in \eqref{CFStage:eq:TPF15} specifies the entire
original inertia representation, including its kernel.
For clarity its image order can also be read off.
Define
\[
 h_k=\begin{cases}1,&N=2\ \hbox{and }k\ \hbox{even},\\
                  N,&\hbox{otherwise}.
       \end{cases}
 \tag{TPF.20}\label{CFStage:eq:TPF20}
\]
The tame image has order \(h_k\). Indeed for \(N=2\)
there is only one character, whose \(k\)th power has
that order. For \(N\ge3,k=1\) the list includes a
primitive character; for \(N\ge3,k\ge2\) every product
character occurs by \eqref{CFStage:eq:TPF18}. In either case
the tame kernel is trivial. When \(w_k\ne0\), the wild
image has order \(p\), and no nontrivial tame scalar
can cancel a nontrivial wild scalar: their orders
divide the coprime integers \(N\) and \(p\).
Thus the full image has order \(h_kp\) in that case,
and order \(h_k\) when \(w_k=0\).

\subsection{Swan conductor and the finite-cover logarithm}

If \(w_k\ne0\), the nontrivial character
\(\mathcal L_\psi(w_k/u)\) has Swan conductor one.
One may verify it directly over \(\bar k_0((u))\):
in \(z^p-z=bw_k/u\), the ramification index is \(p\)
by the valuation argument above, \(v_E(z)=-1\), and
\[
 v_E\bigl((z+a)^{-1}-z^{-1}\bigr)=2
                         \qquad(a\in\mathbf F_p^\times).
 \tag{TPF.21}\label{CFStage:eq:TPF21}
\]
Thus its lower ramification groups are
\(G_0=G_1=\mathbf F_p,G_2=1\), giving one positive
break of value one. Tensoring by a tame character
does not change this wild restriction. Every tuple
line in \eqref{CFStage:eq:TPF9} has this same wild character.
If \(w_k=0\), every tuple line is tame. Additivity of
Swan under direct sum therefore proves
\[
 \boxed{\displaystyle
 \operatorname{Swan}_0(\mathcal T_k)=
 \begin{cases}
 0,&w_k=0,\\
 (N-1)^k,&w_k\ne0 .
 \end{cases}}
 \tag{TPF.22}\label{CFStage:eq:TPF22}
\]
The finite cover
\[
 u=v^N,\qquad z_k^p-z_k=\frac{bkc}{v^N}
 \tag{TPF.23}\label{CFStage:eq:TPF23}
\]
trivializes all geometric inertia of \(\mathcal T_k\).
The first map kills every product Kummer character.
The second trivializes the displayed additive factor
by its own torsor. If \(w_k=0\), it is a trivial
disjoint torsor and the first map alone suffices.
If \(w_k\ne0\), its degree is \(p\), as its pole order
\(N\) is prime to \(p\). When \(c\ne0\) and
\(p\nmid k\), it is explicitly the original cover
\eqref{CFStage:eq:TPF13} with \(z_k=kz\); this equality follows
from \(k^p=k\) in \(\mathbf F_p\).
When \(p\mid k\), the original cover still kills the
representation, and its wild deck subgroup already
acts trivially by \eqref{CFStage:eq:TPF15}.

On the subgroup defined by this finite cover every
monodromy matrix is the identity. Hence its unipotent
logarithm is exactly
\[
 N_{\mathrm{mon},k}=0 .
 \tag{TPF.24}\label{CFStage:eq:TPF24}
\]
There is no conflict with any of the invariant
dimensions in \eqref{CFStage:eq:TPF19}. Those concern the
original inertia group; after the displayed cover,
the restricted representation has the full
\(r_k\)-dimensional invariant space.

\subsection{Frobenius on the invariant space and exact descent}

Let \(j:\mathbf G_m\hookrightarrow\mathbf A^1_u\).
The ordinary invariant stalk
\((j_*\mathcal T_{0,k})_{\bar0}=V_k^{I_0}\):
sections over the punctured strict henselian
neighborhood are precisely the invariant vectors
of its geometric generic stalk. If \(w_k\ne0\)
this stalk is zero. For \(w_k=0\), set
\[
 \Omega_k^0=\{\boldsymbol\chi\in\Omega_k:
                         \eta_{\boldsymbol\chi}=1\}.
 \tag{TPF.25}\label{CFStage:eq:TPF25}
\]
Over \(L\), the exact Frobenius module is
\[
 V_k^{I_0}\simeq
       \bigoplus_{\boldsymbol\chi\in\Omega_k^0}
                    \overline{\mathbf Q}_\ell(A_{\boldsymbol\chi}),
 \qquad
 \det(1-T\operatorname{Fr}_Q\mid V_k^{I_0})
       =\prod_{\boldsymbol\chi\in\Omega_k^0}
                              (1-A_{\boldsymbol\chi}T).
 \tag{TPF.26}\label{CFStage:eq:TPF26}
\]
Each surviving summand is the constant sheaf
\(\mathcal A_{A_{\boldsymbol\chi}}\): its product
Kummer character and its additive factor are
both trivial. This proves the formula for the
boundary stalk, as well as for its constant
extension to the \(u\)-line.

Over \(k_0\), Frobenius need not preserve the
individual tuple lines. Let \(F_0\) denote its
geometric Frobenius action. Its exact permutation
is
\[
 F_0:\boldsymbol\chi\longmapsto
       \boldsymbol\chi^{Q_0}
              =(\chi_1^{Q_0},\ldots,\chi_k^{Q_0}).
 \tag{TPF.27}\label{CFStage:eq:TPF27}
\]
To prove its direction, over the splitting field
the deck operator satisfies
\(F_0D_\zeta F_0^{-1}=D_{\zeta^{Q_0^{-1}}}\);
the exponent denotes the inverse modulo \(N\).
Conjugating the exact projector gives
\[
 F_0\left(\frac1N\sum_\zeta
       \lambda_\chi(\zeta)D_\zeta\right)F_0^{-1}
  =\frac1N\sum_\eta
       \lambda_\chi(\eta^{Q_0})D_\eta
  =e_{\chi^{Q_0}} .
 \tag{TPF.28}\label{CFStage:eq:TPF28}
\]
Tensoring this identity in the retained order proves
\eqref{CFStage:eq:TPF27}. The coefficients of the projectors
are in the fixed \(\ell\)-adic coefficient field;
the conjugation changes their deck operators, not
their coefficient values. The subset \(\Omega_k^0\)
is stable because taking a \(Q_0\)th power preserves
the trivial product character.

For each orbit \(\mathcal O\) in \(\Omega_k^0\),
choose a tuple \(\boldsymbol\chi\) in it and let its
exact orbit length be \(e\). Then \(e\mid d\).
Put \(L_e=\mathbf F_{Q_0^e}\subset L\).
Each \(\chi_i\) uniquely descends by norm to a
nontrivial character
\(\kappa_{i,e}:L_e^\times\to
\overline{\mathbf Q}_\ell^\times\):
\[
 \chi_i=\kappa_{i,e}\circ\operatorname N_{L/L_e},
 \qquad
 \psi_e=\psi_0\circ\operatorname{Tr}_{L_e/k_0}.
 \tag{TPF.29}\label{CFStage:eq:TPF29}
\]
Indeed \(\chi_i^{Q_0^e}=\chi_i\), so its order
divides \(Q_0^e-1\). If \(g\) generates \(L^\times\),
the kernel of this norm is generated by
\(g^{Q_0^e-1}\), on which \(\chi_i\) is trivial.
The norm is surjective, proving existence,
uniqueness, and nontriviality of \(\kappa_{i,e}\).
Also \(\prod_i\kappa_{i,e}=1\), again by surjectivity
of the norm.

It is not necessary that all \(N\)th roots of unity
belong to \(L_e\). Here is the exact source map
giving each individual descent. Put
\(r_i=\operatorname{ord}(\chi_i)\). Then
\(r_i\mid N\) and \(r_i\mid Q_0^e-1\), so
\(\mu_{r_i}\subset L_e\). Factor the original
power map by
\[
 \mathbf A^1_y\longrightarrow\mathbf A^1_x
       \longrightarrow\mathbf A^1_z,\qquad
 x=y^{N/r_i},\qquad z=x^{r_i}=y^N .
\]
Pullback of functions on geometric fibre roots
identifies the direct image for \(x^{r_i}=z\)
with the \(\mu_{N/r_i}\)-invariants of the
original \(y^N=z\) direct image. For \(z\ne0\),
the quotient of its root set by
\(\mu_{N/r_i}\) is exactly the root set of
\(x^{r_i}=z\), and invariant functions are
exactly the functions constant on these
orbits. At zero both geometric fibres have
one point and the map is the identity.
The sheaf map is therefore an isomorphism
onto those invariants at every geometric stalk.

The character \(\lambda_{\chi_i}\) factors
through \(\mu_N\to\mu_{r_i}\),
\(\zeta\mapsto\zeta^{N/r_i}\), as
\(\lambda_{\kappa_{i,e}}\). Indeed for
\(a\in L^\times\),
\[
 \lambda_{\kappa_{i,e}}(a^{(Q-1)/r_i})
   =\kappa_{i,e}(a^{(Q-1)/(Q_0^e-1)})
   =\chi_i(a)
   =\lambda_{\chi_i}(a^{(Q-1)/N}),
\]
and \(a\mapsto a^{(Q-1)/N}\) surjects onto
\(\mu_N\). Thus the descended individual
source summand is precisely the trace-
\(\kappa_{i,e}\) Kummer line, including
its zero stalk. The phase on that source
remains \(z/(Nu)+c/u\), with the original
denominator \(N\), not \(r_i\).
Applying the same source isomorphism and
trace formula as in SPF over \(L_e\)
therefore computes the cycle scalar exactly:
\[
 \Lambda_{\mathcal O}
   =\prod_{i=1}^k\left(
        -\kappa_{i,e}(N)G_{L_e}(\kappa_{i,e},\psi_e)\right).
 \tag{TPF.30}\label{CFStage:eq:TPF30}
\]
Every individual factor, including \(N\) and the minus
sign, is retained. More explicitly \(F_0^e\) acts
by this product on the selected tuple line: at
\(u=1\) all its Kummer values are one, while
the product of its additive values is
\(\psi_e(kc)=1\), since \(w_k=0\).
The tuple product is a constant sheaf, so this
is its action at the invariant boundary stalk too.
This proves the scalar over its actual orbit
field; it has not been selected as an arbitrary
root of an eigenvalue over \(L\).

Choose \(v_0\ne0\) in the tuple line and set
\(v_j=F_0^jv_0\) for \(0\le j<e\).
These vectors are in distinct direct-summand lines,
so form a basis for the orbit block. In this
basis \(F_0v_j=v_{j+1}\) for \(j<e-1\), and
\(F_0v_{e-1}=\Lambda_{\mathcal O}v_0\).
Consequently
\[
 \boxed{\begin{aligned}
 \det(T-F_0\mid V_k^{I_0})
    &=\prod_{\mathcal O\subset\Omega_k^0}
                              (T^{|\mathcal O|}-\Lambda_{\mathcal O}),\\
 \det(1-TF_0\mid V_k^{I_0})
    &=\prod_{\mathcal O\subset\Omega_k^0}
                              (1-\Lambda_{\mathcal O}T^{|\mathcal O|}).
 \end{aligned}}
 \tag{TPF.31}\label{CFStage:eq:TPF31}
\]
These expressions concern \(w_k=0\); when \(w_k\ne0\)
both determinants on the zero space are one.
Choosing another representative does not change
\(\Lambda_{\mathcal O}\), because \(F_0^e\) commutes
with \(F_0\) and hence has the same scalar on every
line of that orbit. Formula \eqref{CFStage:eq:TPF31} thus
proves descent without asserting diagonal
Frobenius in the original tuple coordinates over
\(k_0\).

\subsection{Exact weight and the two cancellation cases}

The elementary unscaled Gauss calculation in SPF
gives, for every complex embedding,
\[
 |G_L(\chi,\psi)|=Q^{1/2},\qquad
 |A_{\boldsymbol\chi}|=Q^{k/2}.
 \tag{TPF.32}\label{CFStage:eq:TPF32}
\]
In \eqref{CFStage:eq:TPF10} all other factors are roots of
unity. At a point of degree \(a\) over \(L\), every
eigenvalue therefore has absolute value \(Q^{ak/2}\).
This proves pointwise purity of weight \(k\) for
the actual ordered tensor. It descends to \(k_0\):
a positive power of any original Frobenius is
covered by the calculation over \(L\), and
\(|\beta^a|=Q_0^{fak/2}\) forces
\(|\beta|=Q_0^{fk/2}\) at an original point of
degree \(f\).

When \(w_k=0\), the same calculation over \(L_e\)
in \eqref{CFStage:eq:TPF30} gives
\[
 |\Lambda_{\mathcal O}|=Q_0^{ek/2},\qquad
 |\beta|=Q_0^{k/2}
       \quad\hbox{for every root }\beta^e=\Lambda_{\mathcal O}.
 \tag{TPF.33}\label{CFStage:eq:TPF33}
\]
Thus every nonzero invariant block has weight
\(k\) on \(\operatorname{Spec}k_0\), including the
blocks that occur because \(p\mid k\).
If \(c=0\), the original single-factor wild
character was already trivial. If \(c\ne0\) and
\(p\mid k\), it becomes trivial by the exact
tensor character \(\vartheta(-ka)\).
Both branches have the same tame multiplicity
formula, but the displayed maps account for
their different origins. No complex arithmetic
zeta zero or comparison of such a zero with these
finite-field Frobenius eigenvalues is constructed
by this calculation.

\subsection{Exact counts of the arithmetic orbit blocks}

For \(f\mid d\), retain the full greatest common
divisor and define
\[
 g_f=\gcd(N,Q_0^f-1),\qquad
 D(f)=\frac{(g_f-1)^k+(g_f-1)(-1)^k}{g_f}.
 \tag{TPF.34}\label{CFStage:eq:TPF34}
\]
The characters fixed by exponent \(Q_0^f\)
form a cyclic subgroup of \(X_N\) of order
\(g_f\). A fixed ordered tuple has all
entries in this subgroup. Applying the
root-of-unity count \eqref{CFStage:eq:TPF18}
inside that subgroup shows that exactly
\(D(f)\) tuples in \(\Omega_k^0\) are fixed.
When \(g_f=1\) there are no allowed
nontrivial entries and the formula is
zero, since \(k\ge1\).
If \(C(e)\) denotes the number of orbits
in \(\Omega_k^0\) of exact length \(e\mid d\), then
\[
 D(f)=\sum_{e\mid f}eC(e),\qquad
 C(e)=\frac1e\sum_{f\mid e}\mu(e/f)D(f).
 \tag{TPF.35}\label{CFStage:eq:TPF35}
\]
The first identity counts the elements
of every orbit fixed by \(F_0^f\).
For the second, multiply that identity
by \(\mu(e/f)\) and sum over \(f\mid e\).
The coefficient of \(rC(r)\) is
\(\sum_{r\mid f\mid e}\mu(e/f)\), which
is one for \(r=e\) and otherwise zero:
the latter sum factors as
\(\prod_{q\mid e/r}(1-1)\) over its
distinct prime divisors. This proves
the orbit count by integer arithmetic.
These orbit numbers count the nonzero boundary
blocks exactly when \(w_k=0\); when \(w_k\ne0\)
the boundary stalk is zero by \eqref{CFStage:eq:TPF19}.

\subsection{The terminal character, its orbit, and the original scalar}

Fix a generator of the abstract cyclic deck group of
order \(N\), and identify its chosen complex \(N\)th
root of unity and its finite-field root in \(\mu_N(L)\)
with a chosen \(N\)th root in the \(\ell\)-adic coefficient
field. These are explicit identifications of finite
cyclic groups. Let \(\chi_*\in X_N\) be the primitive
character whose \(\lambda_{\chi_*}\) is the resulting
tautological character.

The direction of the deck action matters. For
\(\delta_\zeta(y)=\zeta y\), direct pullback gives
\[
 \delta_\zeta^*(y^b\,dy)=\zeta^{b+1}y^b\,dy,\qquad
 \delta_\zeta^*e_r=e_{\zeta^{-1}r},\qquad
 D_\zeta=(\delta_\zeta^{-1})^* .
 \tag{TPF.36}\label{CFStage:eq:TPF36}
\]
The middle equality follows by evaluating the
indicator function of the point \(r\) at \(\zeta y\).
Thus the finite SPF summand labelled \(\chi\) has
\(D_\zeta\)-eigenvalue \(\lambda_\chi(\zeta)^{-1}\)
but ordinary deck-pullback eigenvalue
\(\lambda_\chi(\zeta)\). Comparing the pullback
characters with the pullback characters, rather
than with the inverse action \(D_\zeta\), gives
the exact character-label map
\[
 b\longmapsto\chi_b=\chi_*^{\,b+1},\qquad
 \chi_{\mathrm{ter}}=\chi_*^{\,m}
                  =\chi_*^{-1},\qquad
 \operatorname{ord}(\chi_{\mathrm{ter}})=N .
 \tag{TPF.37}\label{CFStage:eq:TPF37}
\]
Here the terminal coefficient is \(y^{m-1}\), and
its form is \(y^{m-1}dy\), so \(b=m-1=N-2\).
The ordered terminal form has deck-pullback character
\(\zeta^{km}=\zeta^{-k}\), precisely as the product
\(\chi_{\mathrm{ter}}^k\) records.
This proves a map of the finite deck-character
index sets. It does not assert a de Rham--etale
comparison isomorphism or identify a complex
arithmetic operator with finite-field Frobenius.

The repeated terminal tuple is
\(\boldsymbol\chi_{\mathrm{ter}}
=(\chi_{\mathrm{ter}},\ldots,\chi_{\mathrm{ter}})\).
Its actual finite-field summand on the same
\(u\)-line is exactly
\[
 \mathcal L_\psi(kc/u)\otimes
 \mathcal A_{\left[-\chi_{\mathrm{ter}}(N)
                    G_L(\chi_{\mathrm{ter}},\psi)\right]^k}
                  \otimes\mathcal K_{\chi_{\mathrm{ter}}^k}(u).
 \tag{TPF.38}\label{CFStage:eq:TPF38}
\]
Its original-inertia invariants have dimension
\[
 \begin{cases}
 1,&N\mid k\ \hbox{and }kc=0,\\
 0,&\hbox{otherwise}.
 \end{cases}
 \quad
 \begin{array}{ll}
 c=0:&\text{the nonzero case is }N\mid k,\\
 c\ne0:&\text{the nonzero case is }pN\mid k.
 \end{array}
 \tag{TPF.39}\label{CFStage:eq:TPF39}
\]
Indeed its tame product is trivial exactly when
\(N\mid k\), since \(\chi_{\mathrm{ter}}\) is primitive.
Its wild character is trivial exactly when \(kc=0\).
In the second branch \(p\mid k\) and \(N\mid k\)
are equivalent to \(pN\mid k\), since \(\gcd(p,N)=1\).

This single terminal line need not descend to
\(k_0\). Its exact arithmetic orbit consists of
the \(d\) distinct ordered tuples
\[
 \left(\chi_{\mathrm{ter}}^{Q_0^j},\ldots,
             \chi_{\mathrm{ter}}^{Q_0^j}\right),
             \qquad 0\le j<d.
\]
Its length is exactly \(d\): equality of two
tuple entries requires equality of their primitive
characters, hence \(Q_0^e=1\bmod N\), whose least
positive solution is \(d\).
If the criterion in \eqref{CFStage:eq:TPF39} holds, all
these lines are invariant and the descended
terminal-orbit block has dimension \(d\), with
\[
 \begin{split}
 \Lambda_{\mathrm{ter}}
   &=\left[-\chi_{\mathrm{ter}}(N)
                    G_L(\chi_{\mathrm{ter}},\psi)\right]^k,\\
 \det(T-F_0\mid\text{terminal orbit})
   &=T^d-\Lambda_{\mathrm{ter}},\\
 \det(1-TF_0\mid\text{terminal orbit})
   &=1-\Lambda_{\mathrm{ter}}T^d .
 \end{split}
 \tag{TPF.40}\label{CFStage:eq:TPF40}
\]
This is \eqref{CFStage:eq:TPF30}--\eqref{CFStage:eq:TPF31} with
its actual orbit field \(L\). Every root has
absolute value \(Q_0^{k/2}\).
If the criterion fails, the invariant part of
this orbit is zero. No Frobenius eigenline over
\(k_0\) is inferred from a single summand over \(L\).

Finally retain the original scalar multiplying
the terminal coefficient in the full marked
packet \(h_0\), evaluated on a retained stage
where its original source coefficients belong
to the domain of the actual map \(\alpha\):
\[
 K=k(m-1),\qquad
 u_0=(g/h_0)(\rho),\qquad
 C_{h_0,k}=
        \frac{K!}{((m-1)!)^k}\,u_0^k .
 \tag{TPF.41}\label{CFStage:eq:TPF41}
\]
Here \(K\) is an integer, and \(u_0\) is the
specified marked coefficient in that actual
coefficient model. This calculation does not
extend \(\alpha\) to a larger cyclic CRT stage
by assigning values to new denominators.
If such a stage inverts \(p\), no extension
of \(\alpha\) to it exists: applying a purported
extension to \(1=p\,p^{-1}\) would give \(1=0\).
The absence of a stage specialization is not
a vanishing coefficient specialization.
Writing a bar for that same specialization gives
the exact image
\[
 \overline{C}_{h_0,k}
  =\overline{K!}\,
       \bigl(\overline{(m-1)!}\bigr)^{-k}
                           \overline u_0^{\,k},\qquad
 \overline C_{h_0,k}\ne0
    \ \Longleftrightarrow\
           K<p\ \hbox{and }\overline u_0\ne0 .
 \tag{TPF.42}\label{CFStage:eq:TPF42}
\]
The denominator is nonzero because \(m-1=N-2<p\).
The numerator \(K!\) is nonzero modulo \(p\)
exactly when \(K<p\), and a positive power of a
field element is nonzero exactly when that
element is nonzero. These facts prove the
equivalence. When the original unit class and
its inverse are retained at this marked point,
their product is one after evaluation at
\(\rho\) and after specialization; hence
\(\overline u_0\ne0\) in that specified model.
Even there the factorial can vanish.

In particular, when \(m\ge2\) and \(p\mid k\),
one has \(K=k(m-1)\ge p\), so
\(\overline C_{h_0,k}=0\). Thus the \(c\ne0\)
terminal-invariant degrees \(pN\mid k\), on
a marked stage admitting the specified
\(\alpha\), have zero specialized original
scalar for \(m\ge2\),
although their sheaf summand and invariant
orbit block computed above are nonzero.
For \(m=1\), \(K=0\) and the scalar is exactly
\(u_0^k\), so only its original coefficient
specialization can make it zero.
For \(c=0\), the tame condition \(N\mid k\)
and the separate scalar condition
\(k(m-1)<p,\ \overline u_0\ne0\) are both
retained.

The last scalar calculation has a precise
coefficient-module target. The original
translation and coefficient specialization give
\[
 \begin{split}
 R[y_1,\ldots,y_k]/(y_1^m,\ldots,y_k^m)
    &\longrightarrow
 \overline E_k
   =k_0[y_1,\ldots,y_k]/(y_1^m,\ldots,y_k^m),\\
 \sum_{\boldsymbol a}q_{\boldsymbol a}
                         \prod_i y_i^{a_i}
    &\longmapsto
       \sum_{\boldsymbol a}\alpha(q_{\boldsymbol a})
                         \prod_i y_i^{a_i},\\
 \overline\tau_{\mathrm{ter}}&=\prod_i y_i^{m-1}\ne0,\qquad
 \overline C_{h_0,k}\,\overline\tau_{\mathrm{ter}}=0
       \ \Longleftrightarrow\ \overline C_{h_0,k}=0 .
 \end{split}
 \tag{TPF.43}\label{CFStage:eq:TPF43}
\]
Successive division by the monic polynomials
\(y_i^m\) gives the unique remainder basis
\(\prod_i y_i^{a_i}\), \(0\le a_i<m\), so
\(\overline\tau_{\mathrm{ter}}\) is one of its nonzero basis
vectors and the last equivalence follows.
All nilpotent directions of this coefficient
algebra are retained. These scalar statements
are in characteristic \(p\). There is no
unital field map from \(k_0\) to
\(\overline{\mathbf Q}_\ell\), since such a
map would send \(p\cdot1=0\) to the nonzero
element \(p\cdot1\). Thus no multiplication
of an \(\ell\)-adic Gauss line by
\(\overline C_{h_0,k}\) is being asserted.
The deck-character map \eqref{CFStage:eq:TPF37}
does not by itself transport coefficient
vectors to etale cohomology; the
finite-field sheaf is the independently
specified compact-support sheaf and does
not depend on this scalar.

\subsection{The shared integral deck-character model and its two base changes}

There is an exact common integral representation
behind the character correspondence, which
does not require a map between fields of
different characteristics. Let
\[
 \begin{split}
 B&=\mathbf Z[\zeta_N,1/N],\qquad
 C_N=\langle\gamma:\gamma^N=1\rangle,\\
 e_j&=\frac1N\sum_{a=0}^{N-1}
                   \zeta_N^{-ja}\gamma^a
             \in B[C_N]\quad(0\le j<N).
 \end{split}
 \tag{TPF.44}\label{CFStage:eq:TPF44}
\]
Here \(\zeta_N\) is a fixed primitive root in the
cyclotomic ring. Finite geometric sums give
\[
 e_ie_j=\delta_{ij}e_j,\qquad
 \sum_{j=0}^{N-1}e_j=1,\qquad
 \gamma e_j=\zeta_N^j e_j .
 \tag{TPF.45}\label{CFStage:eq:TPF45}
\]
For example the coefficient convolution in
\(e_ie_j\) contains
\(\sum_a\zeta_N^{a(j-i)}\), which is \(N\)
for \(i=j\) and zero otherwise. The sum of
the \(e_j\) has coefficient one at
\(\gamma^0\) and zero elsewhere, and
reindexing \(a+1\) proves the last identity.

The corresponding free module of polynomial
forms with the original degree retained is
\[
 M_B=B[y]\,dy/y^mB[y]\,dy,\qquad
 \gamma(y^bdy)=\zeta_N^{b+1}y^bdy
                      \quad(0\le b<m).
 \tag{TPF.46}\label{CFStage:eq:TPF46}
\]
It is the quotient of polynomial one-forms
by \(y^m\) times that free module, not the
module of Kahler differentials of the
truncated algebra. Its basis is
\(dy,y\,dy,\ldots,y^{m-1}dy\).
Thus the exact \(B[C_N]\)-linear isomorphism
\[
 \bigoplus_{j=1}^{N-1}B e_j
       \longrightarrow M_B,\qquad
 e_j\longmapsto y^{j-1}dy
 \tag{TPF.47}\label{CFStage:eq:TPF47}
\]
follows from \eqref{CFStage:eq:TPF45} and the displayed
bases. In an ordered \(k\)-fold tensor,
\(e_{j_1}\otimes\cdots\otimes e_{j_k}\)
therefore projects onto the corresponding
ordered monomial-form basis line. The
simultaneous \(e_0\)-projector selects
\(\sum_i j_i=0\bmod N\), exactly the count
in \eqref{CFStage:eq:TPF18}. The terminal line is
\((B e_m)^{\otimes k}\).

Choose a primitive root \(\zeta_L\in L\)
and the primitive \(\ell\)-adic root
\(\zeta_\ell\) already used in
\eqref{CFStage:eq:TPF37}. The two actual ring maps are
\[
 \iota_p:B\to L,\quad\zeta_N\mapsto\zeta_L,
 \qquad
 \iota_\ell:B\to\overline{\mathbf Q}_\ell,\quad
                  \zeta_N\mapsto\zeta_\ell .
 \tag{TPF.48}\label{CFStage:eq:TPF48}
\]
They are well defined: \(p\nmid N\),
\(N\mid Q-1\), and both selected roots are
primitive \(N\)th roots. The first base
change of \eqref{CFStage:eq:TPF47} is the concrete
isomorphism with
\(L[y]dy/y^mL[y]dy\), under which
\(e_j\) is the projector onto
\(L y^{j-1}dy\).
The coefficient-map connection is also exact:
\(R\otimes_{\mathbf Z}B\to L\) sends
\(r\otimes b'\) to \(\alpha(r)\iota_p(b')\).
Tensoring its polynomial-form module gives
this same characteristic-\(p\) module.
This does not introduce new inverses from
a larger marked stage.

For the second base change, let \(W_\ell\)
be the fibre at \(u=1\) of the actual
untwisted sheaf
\(\mathcal L_\psi(c/u)^\vee\otimes\mathcal F\).
Give it the source deck-pullback action
\(\gamma=D_{\zeta_L}^{-1}\).
The image of the same integral projector is
\[
 \iota_\ell(e_j)
   =\frac1N\sum_{a=0}^{N-1}
        \zeta_\ell^{-ja}D_{\zeta_L}^{-a}
   =\frac1N\sum_{a=0}^{N-1}
        \zeta_\ell^{ja}D_{\zeta_L}^{a}
   =e_{\chi_*^j}^{\mathrm{SPF}} .
 \tag{TPF.49}\label{CFStage:eq:TPF49}
\]
The second equality reindexes \(a\mapsto-a\).
The third is precisely the original SPF
power-map projector, with
\(\lambda_{\chi_*^j}(\zeta_L)=\zeta_\ell^j\).
Thus it projects onto the actual
one-dimensional Gauss line
\(H_j=e_{\chi_*^j}^{\mathrm{SPF}}W_\ell\),
whose geometric Frobenius over \(L\) is
\[
 -\chi_*^j(N)G_L(\chi_*^j,\psi).
\]
Keep that entire line, including its
Frobenius, rather than choosing a numerical
basis for it. If \(H_j^{\mathrm{triv}}\)
denotes the same vector space and Frobenius
with only its \(C_N\)-action replaced by
the trivial action, there is the canonical
isotypic reconstruction below. Give each
first character-module factor the identity
splitting-field Frobenius:
\[
 \begin{split}
 \bigoplus_{j=1}^{N-1}
    \left(B e_j\otimes_{B,\iota_\ell}
                     \overline{\mathbf Q}_\ell\right)
            \otimes_{\overline{\mathbf Q}_\ell}
                         H_j^{\mathrm{triv}}
       &\longrightarrow W_\ell,\\
 (a e_j)\otimes h&\longmapsto a h .
 \end{split}
 \tag{TPF.50}\label{CFStage:eq:TPF50}
\]
It is an isomorphism: on the \(j\)th
summand it is the scalar multiplication
isomorphism onto \(H_j\), and the SPF
projectors are complete and orthogonal.
It preserves the deck action by
\eqref{CFStage:eq:TPF45} and preserves the
splitting-field Frobenius by the retained
action on \(H_j^{\mathrm{triv}}\).
No Frobenius action on a characteristic-\(p\)
coefficient vector has been assumed.

Tensoring these actual maps in the fixed
order gives the shared terminal-character
projector on both sides:
\((e_m)^{\otimes k}\) maps to the
characteristic-\(p\) line
\(L(\prod_i y_i^{m-1}dy_i)\) and to
the ordered Gauss line \(H_m^{\otimes k}\).
The full finite-field sheaf reinstates
its already specified additive factor
\(\mathcal L_\psi(kc/u)\); the common
integral model supplies the finite
\(\mu_N\) representation, not an
unconstructed universal additive sheaf.
Formulas \eqref{CFStage:eq:TPF44}--\eqref{CFStage:eq:TPF50}
are the exact common projector and
base-change correspondence. They retain
the coefficient scalar specialization
in \eqref{CFStage:eq:TPF43} and the independently
nonzero \(\ell\)-adic sheaf lines as
their stated, different targets.

\subsection{The multiplicity-one case and its exact scalar comparison}

When \(m=1\), one has \(N=2\), \(p>2\),
\(d=1\), and \(L=k_0=\mathbf F_Q\), because
every odd \(Q_0\) satisfies \(Q_0=1\bmod2\).
There is just the nontrivial quadratic
character \(\chi\), and it is the terminal
character. Retain the original constants
\[
 c=-\rho^2/2,\qquad
 \gamma_\chi=-\chi(2)G_L(\chi,\psi).
 \tag{TPF.51}\label{CFStage:eq:TPF51}
\]
Their exact Gauss-square identity is
\[
 \gamma_\chi^2=\chi(-1)Q.
 \tag{TPF.52}\label{CFStage:eq:TPF52}
\]
To prove it, take a complex embedding of the
cyclotomic field of character values.
Quadratic-character values are real, so
changing \(z\) to \(-z\) in the conjugate
Gauss sum gives
\(\overline{G_L(\chi,\psi)}
 =\chi(-1)G_L(\chi,\psi)\).
The proved absolute-square identity
\(G_L\overline{G_L}=Q\) implies
\(G_L^2=\chi(-1)Q\), since \(\chi(-1)^2=1\).
Also \(\chi(2)^2=1\), proving
\eqref{CFStage:eq:TPF52} with the original minus
sign and the factor \(2\) retained.
Both sides lie in the same cyclotomic
number field; injectivity of the chosen
embedding proves that algebraic identity
there and hence in its chosen \(\ell\)-adic
realization.

For \(c\ne0\), every positive terminal-invariant
degree is exactly \(k=2pr\), \(r\ge1\).
Its one-dimensional original-base invariant
Frobenius is
\[
 \gamma_\chi^k
      =(\chi(-1)Q)^{pr},\qquad
 |\gamma_\chi^k|=Q^{k/2}.
 \tag{TPF.53}\label{CFStage:eq:TPF53}
\]
Its additive factor is exactly
\(\mathcal L_\psi(2prc/u)=\mathcal L_\psi(0)\);
its Kummer factor is
\(\mathcal K_{\chi^{2pr}}=\mathcal K_1\).
Thus the line is an actual nonzero invariant
line of weight \(k\).
Here \(K=k(m-1)=0\) and the original marked
coefficient is exactly \(u_0^k\).
On an actual coefficient stage retaining
the inverse of \(u_0\), its reduction is
nonzero. There is no factorial loss.
For \(c=0\), every even \(k\) is tame
invariant and the same formula reads
\(\gamma_\chi^k=(\chi(-1)Q)^{k/2}\);
no divisibility by \(p\) is needed.

There is also an exact comparison calculation
for any chosen map between the corresponding
one-dimensional complex spaces, without
asserting a canonical comparison.
The actual arithmetic primary algebra is
\(\mathbf C[s]/(s-\rho)\). In its ordered
\(k\)-fold tensor, the original sum operator
\(S=\sum_i s_i\) is multiplication by \(k\rho\).
For \(a>0\), its dilation operator is therefore
\[
 A_a=\exp((\log a)S)=a^{k\rho}\operatorname{id},
 \qquad a^{k\rho}=\exp(k\rho\log a).
 \tag{TPF.54}\label{CFStage:eq:TPF54}
\]
Choose a complex embedding of the cyclotomic
number field containing \(\gamma_\chi\),
and let \(L_{\mathrm{G}}\) be the resulting
one-dimensional complex Frobenius model,
with \(F_{\mathrm G}=\gamma_\chi^k\operatorname{id}\).
This takes a number-field model of the
computed eigenvalue to \(\mathbf C\), not
a field map from \(k_0\) to \(\mathbf C\).
For every nonzero complex-linear map
\(T:L_{\mathrm G}\to L_{\mathrm{arith}}\),
scalar multiplication proves
\[
 \begin{split}
 A_aT
    &=\frac{a^{k\rho}}{\gamma_\chi^k}\,TF_{\mathrm G},\\
 A_aT-TF_{\mathrm G}
    &=(a^{k\rho}-\gamma_\chi^k)T,\qquad
 \left|\frac{a^{k\rho}}{\gamma_\chi^k}\right|
       =\frac{a^{k\operatorname{Re}\rho}}{Q^{k/2}} .
 \end{split}
 \tag{TPF.55}\label{CFStage:eq:TPF55}
\]
For example such maps exist by sending a
nonzero basis vector to a nonzero basis
vector; changing either basis multiplies
the matrix of \(T\) by a nonzero scalar
and leaves the displayed ratio unchanged.
At \(a=Q\) the ratio has modulus
\(Q^{k(\operatorname{Re}\rho-1/2)}\).
This is the exact intertwining defect for
every such map. Existence of a nonzero
map does not turn this computed ratio
into one, and no identification of the
two operators has been assumed.

\paragraph{Sources and dependencies.}
The sole specialized input is the proved SPF.1--29
single-primary sheaf calculation, with its exact
power-map projectors, cohomology, Gauss signs,
and local inertia. The foundational curve and
trace formulas are in Katz,
\emph{Gauss Sums, Kloosterman Sums, and Monodromy
Groups}, Sections 2.0--2.3 and 4.0--4.3.
The torsor constructions are the Stacks Project,
tags \texttt{03PK} and \texttt{0A3J}. The relative
compact-support Kunneth map in
\eqref{CFStage:eq:TPF6} is SGA~4, Expose XVII,
Theorem 5.4.3, applied to the common base
\(\mathbf G_{m,u}\). All
tensor multiplicities, original inertia
projectors, cancellation cases, descent
permutations, and orbit scalars are proved here.

% END TPF
% BEGIN TCL SHA256 7371f26aa05099d14d72d55999c9af8605125a20da017876b3e129775e2efaae
\section{The retained primary cyclic lattice at its specified coefficient prime}
\label{CFStage:sec:tcl}

This calculation uses the ordered primary tensor and full Taylor unit
of TP.1 and TP.23--25. It computes its original cyclic map over the
specified coefficient prime, including the map before reduction, its
primitive containing lattice, its cokernel, and the exact reduction
sequence. The support carrier used below is the original
\(G_{C_n}\) of MFC.29 and MFC.35--44.

\subsection{The actual coefficient localization and the unchanged coordinates}

Use the single-primary coefficient specialization already specified
in SPF.1. Denote its retained coefficient subring of \(\mathbb C\)
by \(A\), its map by \(\alpha:A\to k_0\), and its characteristic
by the prime \(p\). Retain every coefficient of the full Taylor
class and its inverse, and put
\[
 \mathfrak q=\ker\alpha,\qquad S=A_{\mathfrak q},\qquad
 \widetilde\alpha:S\longrightarrow k_0,\qquad
 \widetilde\alpha(a/b)=\alpha(a)/\alpha(b).
 \tag{TCL.1}\label{CFStage:eq:tcl-ring}
\]
Thus \(b\notin\mathfrak q\) in this formula. The image of
\(\alpha\) is a finite subring of a field. For each nonzero element
of that image, multiplication on the finite image is injective and
therefore surjective; its inverse lies in the image. The image is
consequently a field, proving that \(\mathfrak q\) is maximal.
The fraction formula is independent of representatives: equality
of two fractions gives equality after cross multiplication in
the domain \(A\), and applying \(\alpha\) and dividing by the
nonzero denominator images gives the asserted equality.
It preserves the ring operations by the same calculation.
The map extends the specified map, rather than selecting a new one.

Since \(A\subset\mathbb C\), both \(A\) and its localization
\(S\) are domains of characteristic zero. The element \(p\) is
a nonzero nonunit: its image is zero whereas the image of the unit
is one. Every integer prime to \(p\) is outside \(\mathfrak q\)
and is a unit in \(S\). In particular \(\mathbb Z_{(p)}\to S\)
is a well-defined injective scalar map. No assertion that \(S\)
is a discrete valuation ring, a principal ideal domain, or a ring
with all ideals in a chain is made or used. The algebra below uses
\(p>m\); the retained SPF specialization has the stronger
\(p>m+1\), as required there by its phase denominators.

Keep the selected primary coordinate \(\rho\), its full order
\(m\geq1\), and the ordered number of factors \(k\geq1\). Put
\[
 \begin{aligned}
 K&=k(m-1),\qquad y_i=s_i-\rho,\qquad
 E_S=S[s_1,\ldots,s_k]/((s_i-\rho)^m)_{i=1}^k
       =S[y_1,\ldots,y_k]/(y_i^m)_{i=1}^k,\\
 J&=\sum_{i=1}^k y_i,\qquad
 M=\sum_{i=1}^k s_i=k\rho+J,\qquad
 \mathcal C_S=S[T]/(T^{K+1}),\qquad
 \mathcal C_S^{\rm orig}=S[Z]/((Z-k\rho)^{K+1}).
 \end{aligned}
 \tag{TCL.2}\label{CFStage:eq:tcl-objects}
\]
The two displayed presentations of \(E_S\) are identified by the
mutually inverse substitutions \(s_i=y_i+\rho\) and
\(y_i=s_i-\rho\). The cyclic-source coordinate map
\(\mathcal C_S^{\rm orig}\to\mathcal C_S\) is
\(Z\mapsto k\rho+T\), with inverse \(T\mapsto Z-k\rho\).
These substitutions respect the defining ideals and their composites
are the identity on each generator. They preserve the original objects
and will be recorded explicitly in the original monomial matrices.

Let
\(\mathcal B=\{0,\ldots,m-1\}^k\), ordered lexicographically.
Repeated division by the monic relations \(y_i^m\) proves that
\(\{y^{\mathbf b}:\mathbf b\in\mathcal B\}\) is an \(S\)-basis:
deleting monomials with any exponent at least \(m\) gives a
remainder, and the relation ideal is spanned by precisely the deleted
monomials, so no relation between retained monomials is possible.
Every surviving degree is at most \(K\); hence \(J^{K+1}=0\).

Write the retained one-factor unit and its retained inverse as
\[
 \begin{split}
 v(y)&=\sum_{j=0}^{m-1}a_jy^j,\qquad
 a_j=\frac{(g/h)^{(j)}(\rho)}{j!}
     =\frac{g^{(m+j)}(\rho)}{(m+j)!},\qquad h=(s-\rho)^m,\\
 w(y)&=\sum_{j=0}^{m-1}b_jy^j,\qquad
 b_0=a_0^{-1},\qquad
 b_r=-a_0^{-1}\sum_{j=1}^{r}a_jb_{r-j},\\
 U&=\prod_{i=1}^kv(y_i),\qquad
 U^{-1}=\prod_{i=1}^kw(y_i)\quad\hbox{in }E_S.
 \end{split}
 \tag{TCL.3}\label{CFStage:eq:tcl-full-unit}
\]
These are the actual retained coefficients, all lying in \(A\)
by its specified construction. The derivative identity follows by
expanding \(g=(s-\rho)^m(g/h)\) at \(\rho\). The recurrence makes
the constant coefficient of \(v w\) equal to one and its coefficients
of degrees \(1,\ldots,m-1\) equal to zero. Thus \(v w=1\) modulo
\(y^m\), proving the inverse in every factor and then in the tensor.
In particular multiplication by \(U\) is an actual \(S\)-linear
automorphism of \(E_S\) commuting with multiplication by \(J\)
and \(M\). No coefficient of this unit has been discarded.

The unchanged unweighted cyclic ring map and the original unit-weighted
cyclic coefficient map are respectively
\[
 \begin{split}
 f:\mathcal C_S\longrightarrow E_S,&\qquad f(T^r)=J^r,\\
 f_U:\mathcal C_S\longrightarrow E_S,&\qquad f_U(T^r)=U J^r,\\
 f_U^{\rm orig}:\mathcal C_S^{\rm orig}\longrightarrow E_S,
 &\qquad f_U^{\rm orig}(Z^j)=U M^j.
 \end{split}
 \tag{TCL.4}\label{CFStage:eq:tcl-original-map}
\]
The relation \(J^{K+1}=0\) proves that all maps are well-defined.
The last two are linear maps, and are maps of modules over the cyclic
source ring acting on the target by \(T\mapsto J\), respectively
\(Z\mapsto M\). They are not asserted to be unital ring maps:
their value at one is the retained \(U\). Their exact intertwining
identities are
\[
 f_U(Tc)=Jf_U(c),\qquad
 f_U^{\rm orig}(Zc)=M f_U^{\rm orig}(c).
 \tag{TCL.5}\label{CFStage:eq:tcl-intertwining}
\]
Both follow by multiplication in the commutative algebra \(E_S\).
This is the coefficient realization of the original source map
\(\mathcal J_k\mathcal T_k\) in TP.25, restricted to polynomials
in the sum of its unchanged \(s_i\) coordinates.

\subsection{The primitive lattice and the exact factorial inclusion}

For \(0\leq r\leq K\), define
\[
 \mathcal B_r=\{\mathbf b\in\mathcal B:|\mathbf b|=r\},\qquad
 \mathbf b!=\prod_{i=1}^k b_i!,\qquad
 V_r=\sum_{\mathbf b\in\mathcal B_r}\frac{y^{\mathbf b}}{\mathbf b!},
 \qquad V_{K+1}=0.
 \tag{TCL.6}\label{CFStage:eq:tcl-divided}
\]
Every denominator is a unit because \(b_i<m<p\). Each set
\(\mathcal B_r\) is nonempty. For \(m>1\) a specified member
\(\mathbf b^*(r)\) is obtained recursively by
\(b_i^*(r)=\min(m-1,r-\sum_{j<i}b_j^*(r))\).
The remaining degree is nonnegative at each step and is zero after
\(k\) steps because \(r\leq k(m-1)\). For \(m=1\) there is only
\(r=0\), with \(\mathbf b^*(0)=(0,\ldots,0)\).

The multinomial expansion and direct multiplication give
\[
 J^r=r!V_r\quad(0\leq r\leq K),\qquad
 J V_r=(r+1)V_{r+1}\quad(0\leq r\leq K).
 \tag{TCL.7}\label{CFStage:eq:tcl-factorial}
\]
To prove the first formula, expand the product of \(r\) copies
of \(\sum_i y_i\). For a prescribed exponent vector \(\mathbf b\)
there are \(r!/\mathbf b!\) assignments of its \(r\) ordered
factors to the variables. Exponents at least \(m\) give zero in
the quotient. This leaves exactly the displayed coefficients.
For the second formula, the coefficient of a retained
\(y^{\mathbf b}\) of degree \(r+1\) is
\(\sum_{i:b_i>0}1/(b_i-1)!\prod_{j\ne i}1/b_j!
=(\sum_i b_i)/\mathbf b!=(r+1)/\mathbf b!\).
For \(r=K\) every resulting monomial is zero, which agrees with
the convention for \(V_{K+1}\).

Define coefficient functionals, a containing lattice and its projection by
\[
 \begin{split}
 \lambda_r(x)&=\mathbf b^*(r)!\,[y^{\mathbf b^*(r)}]x,\\
 \Lambda&=\bigoplus_{r=0}^K S V_r\subset E_S,\qquad
 \pi(x)=\sum_{r=0}^K\lambda_r(x)V_r,\\
 F&=\operatorname{span}_S
       \{y^{\mathbf b}:\mathbf b\notin\{\mathbf b^*(0),\ldots,
                                                   \mathbf b^*(K)\}\}.
 \end{split}
 \tag{TCL.8}\label{CFStage:eq:tcl-splitting}
\]
Distinct \(V_r\) have distinct degrees, and the pivot coefficient
of \(V_r\) is the unit \(1/\mathbf b^*(r)!\). Therefore
\(\lambda_r(V_t)=\delta_{rt}\), proving both the direct sum and
\(\pi^2=\pi\). Subtracting \(\pi(x)\) deletes exactly all pivot
coordinates, so \(x-\pi(x)\in F\). Conversely all pivot
coordinates vanish on \(F\); hence \(\pi(F)=0\). We have proved
the explicit splitting
\[
 E_S=\Lambda\oplus F,\qquad \ker\pi=F.
 \tag{TCL.9}\label{CFStage:eq:tcl-direct}
\]
In particular each \(V_r\) is primitive in the sense that it is
part of the displayed basis of a direct summand, with no division
by \(r!\) in the lattice construction.

Write \(L=f(\mathcal C_S)\). Equations (TCL.7)--(TCL.9) imply
\[
 \begin{split}
 L&=\bigoplus_{r=0}^K r!S V_r\subset\Lambda,\qquad
 \Lambda/L\simeq\bigoplus_{r=0}^K S/(r!),\\
 W_r&=U V_r,\quad \Lambda_U=U\Lambda,
 \quad L_U=f_U(\mathcal C_S)=UL,\quad F_U=UF,\\
 E_S&=\Lambda_U\oplus F_U,\qquad
 \pi_U=U\pi U^{-1},\qquad
 L_U=\bigoplus_{r=0}^K r!S W_r.
 \end{split}
 \tag{TCL.10}\label{CFStage:eq:tcl-unit-lattices}
\]
Since \(S\) is a characteristic-zero domain, each \(r!\) is
nonzero and multiplication by it is injective. Thus \(f\) and
\(f_U\) are injective. The quotient isomorphism sends
\(\sum_r c_r V_r\) to \((c_r\bmod r!)_r\); its kernel and
surjectivity follow coordinate by coordinate. Applying the
automorphism \(U\) proves every transported identity in the
display. There is no claim that \(U\Lambda=\Lambda\); it is the
pair of lattices \(L\subset\Lambda\) that is transported.

Let \(Q=\operatorname{Frac}(S)\). The exact saturation statement is
\[
 (L_U\otimes_S Q)\cap E_S=\Lambda_U
       \quad\hbox{inside }E_S\otimes_S Q.
 \tag{TCL.11}\label{CFStage:eq:tcl-saturation}
\]
Indeed the nonzero \(r!\) are invertible in \(Q\), so
\(L_U\otimes Q=\Lambda_U\otimes Q\). In the splitting
\(E_S=\Lambda_U\oplus F_U\), an integral element belonging
to this rational subspace has zero \(F_U\)-component, because
the free module \(F_U\) injects into \(F_U\otimes Q\).
This proves the stated intersection, including both inclusions.

Both lattices are stable under \(J\). On their indicated bases,
the full action and its covariant inclusion are
\[
 \begin{split}
 J W_r&=(r+1)W_{r+1},\qquad
 J(r!W_r)=(r+1)!W_{r+1},\\
 \mathcal C_S\xrightarrow{\ f_U\ }\Lambda_U,
 &\qquad T^r\longmapsto r!W_r,\qquad
 f_U\circ M_T=M_J\circ f_U.
 \end{split}
 \tag{TCL.12}\label{CFStage:eq:tcl-action-matrix}
\]
At \(r=K\) both actions are zero. Commutation of \(U\) and
\(J\) and (TCL.7) prove all entries. This displays the actual
source multiplication and the containing-lattice multiplication
together with their exact intertwining map.

For clarity, the same map in all original monomial coordinates is
\[
 \begin{split}
 P_{ab}&=\begin{cases}\binom ba\rho^{b-a}&a\leq b,\\0&a>b,
                 \end{cases}\quad(0\leq a,b<m),\qquad P_k=P^{\otimes k},\\
 B_{rj}&=\begin{cases}\binom jr(k\rho)^{j-r}&r\leq j,\\0&r>j,
                 \end{cases}\quad(0\leq r,j\leq K),\\
 \mathsf V_{\mathbf b r}&=
       \begin{cases}1/\mathbf b!&|\mathbf b|=r,\\0&\text{otherwise},
       \end{cases}\qquad \Delta=\operatorname{diag}(0!,1!,\ldots,K!),\\
 [f_U^{\rm orig}]_{(s^{\mathbf b}),(Z^j)}
     &=P_k^{-1}\,[U]_{(y^{\mathbf b})}\,\mathsf V\,\Delta\,B.
 \end{split}
 \tag{TCL.13}\label{CFStage:eq:tcl-original-matrix}
\]
The binomial theorem gives the two translation matrices \(P_k\)
and \(B\). Their inverses replace \(\rho\) and \(k\rho\) by
their negatives, respectively, as verified by the inverse substitutions
above; their diagonal entries are one. The \(j\)-th column on the
right is consequently the unchanged polynomial
\(U\sum_{r=0}^j\binom jr(k\rho)^{j-r}r!V_r=U(k\rho+J)^j\)
written in the original \(s^{\mathbf b}\)-basis. This proves the
matrix formula without replacing the original source map or its unit.

\subsection{Exact integer depths, quotient coordinates, and reduction}

For each integer \(r\geq0\) define
\[
 d_r=v_p(r!)=\sum_{a\geq1}\left\lfloor\frac{r}{p^a}\right\rfloor,
 \qquad r!=p^{d_r}\varepsilon_r,\qquad
 \varepsilon_r\in\mathbb Z_{(p)}^\times\subset S^\times.
 \tag{TCL.14}\label{CFStage:eq:tcl-depth}
\]
Here \(0!=1\), and the sum is finite because its terms vanish
when \(p^a>r\). Each integer \(j\) between one and \(r\)
contributes one factor of \(p\) for each positive \(a\) with
\(p^a\mid j\). Interchanging these two finite counts yields
\(\sum_a\#\{1\leq j\leq r:p^a\mid j\}
=\sum_a\lfloor r/p^a\rfloor\), proving the valuation formula.
The remaining integer factor is prime to \(p\), hence is a unit
in \(S\). Therefore \((r!)=p^{d_r}S\) exactly. The depth cannot
increase in \(S\): an equality \(p^d=p^{d+1}s\) would permit
cancellation in the domain and give \(1=ps\), contradicting
\(\widetilde\alpha(1)=1\) and \(\widetilde\alpha(p)=0\).
This argument proves strictness of every consecutive pair of
principal \(p\)-power ideals, without assigning a valuation to
arbitrary elements of \(S\).

Set \(C=\Lambda_U/L_U\). The following are actual quotient maps:
\[
 \begin{split}
 q_r:E_S&\longrightarrow S/(p^{d_r}),\qquad
 q_r(x)=\lambda_r(U^{-1}x)\bmod p^{d_r},\\
 C&\xrightarrow{\ \sim\ }\bigoplus_{r=0}^K S/(p^{d_r}),
       \qquad [\sum_r c_rW_r]\longmapsto
                       (c_r\bmod p^{d_r})_r,\\
 E_S/L_U&\xrightarrow{\ \sim\ }
                 F_U\oplus\bigoplus_{r=0}^K S/(p^{d_r}).
 \end{split}
 \tag{TCL.15}\label{CFStage:eq:tcl-quotient}
\]
The convention is \(S/(p^0)=0\). Each \(q_r\) is linear,
vanishes on \(F_U\) and on \(L_U\), and takes \(W_t\) to
\(\delta_{rt}\) in its indicated target. Equations
(TCL.9)--(TCL.10) prove the two isomorphisms and all these claims.
For the coordinate class \(\nu_r=[W_r]\in C\), its exact
annihilator in \(S\) is \(p^{d_r}S\), including the unit ideal
when \(\nu_r=0\) and \(d_r=0\).

Let \(B\) denote either \(S/pS\) with its quotient map, or
the specified target \(k_0\) with the map
\(\widetilde\alpha\). In both cases \(B\ne0\) and has
characteristic \(p\). Append a bar to each base-changed module,
element and map. Equations (TCL.8)--(TCL.10) base change directly,
because their projections and inverse changes of basis have their
coefficients in \(S\). Thus \(\overline W_0,\ldots,
\overline W_K\) remain a basis of the direct summand
\(\overline\Lambda_U\subset\overline E\). The original map becomes
\[
 \begin{gathered}
 \overline f_U(\overline T^r)=
 \begin{cases}
  \overline{r!}\,\overline W_r&0\leq r<p,\\
  0&p\leq r\leq K,
 \end{cases}\\
 \begin{aligned}
 \ker\overline f_U
    &=\bigoplus_{r=p}^K B\overline T^r=(\overline T^p),\\
 \operatorname{im}\overline f_U
    &=\bigoplus_{r=0}^{\min(K,p-1)}B\overline W_r.
 \end{aligned}
 \end{gathered}
 \tag{TCL.16}\label{CFStage:eq:tcl-reduced-map}
\]
An empty direct sum is zero, and the ideal \((\overline T^p)\)
is zero when \(p>K\). For \(r<p\), every factor of \(r!\)
is a nonzero element of the prime field and therefore a unit in
\(B\). For \(r\geq p\) the factorial has a factor \(p\)
and maps to zero. Applying these facts to the independent
\(\overline W_r\) proves exactly the displayed kernel and image.

The nilpotency indices of the indicated actions are
\[
 \begin{array}{c|c}
 \text{module and action}&\text{nilpotency index}\\ \hline
 \overline{\mathcal C}_S=B[T]/(T^{K+1}),\ M_T&K+1\\
 \operatorname{im}\overline f_U,\ M_{\overline J}
           &\min(p,K+1)\\
 \overline E,\ M_{\overline J}&\min(p,K+1)\\
 \overline\Lambda_U,\ M_{\overline J}&\min(p,K+1).
 \end{array}
 \tag{TCL.17}\label{CFStage:eq:tcl-indices}
\]
The nilpotency index is the least positive integer whose power is
zero; thus it is one for the zero action on a nonzero module.
In the first row, \(T^K\ne0\) in its free monomial basis and
\(T^{K+1}=0\). On \(\overline E\), the binomial coefficients
\(\binom pi\) for \(0<i<p\) are divisible by \(p\): their
factorial expression has one factor of \(p\) in its numerator
and none in its denominator. The commuting-variable binomial
formula therefore gives
\(\overline J^p=\sum_i\overline y_i^p=0\), since \(p>m\).
Also \(\overline J^{K+1}=0\). If
\(d=\min(p-1,K)\), then
\(\overline J^d=\overline{d!}\,\overline V_d\ne0\), by its
unit pivot coordinate. This proves the third row. Its application
to \(\overline W_0=\overline U\) gives
\(\overline J^d\overline W_0
=\overline{d!}\,\overline W_d\ne0\), proving the second and
fourth rows together with their inherited upper bounds.
In particular the abstract reduced cyclic source has not acquired
the smaller index: its exact map to the target has lost the
directions in (TCL.16).

The reduced containing lattice has the more precise chain decomposition
\[
 \overline\Lambda_U=
 \bigoplus_{\substack{j\geq0\\jp\leq K}}
       \bigoplus_{t=0}^{\min(p-1,K-jp)} B e_{j,t},\qquad
 e_{j,t}=\left(\prod_{a=1}^t\overline{jp+a}\right)
                          \overline W_{jp+t}.
 \tag{TCL.18}\label{CFStage:eq:tcl-divided-chains}
\]
The empty product is one. All displayed products are units because
\(1\leq a\leq p-1\). These are therefore changes of basis on
disjoint intervals of the basis \(\overline W_r\). Formula
(TCL.12) gives \(\overline J e_{j,t}=e_{j,t+1}\) except at the
last position of a displayed interval. There it is zero, either
because the next coefficient is \(\overline{(j+1)p}=0\), or
because the next index is \(K+1\). This proves each full chain,
including every chain that is absent from the image in (TCL.16).

The lost directions are also the exact connecting image in the
base-change sequence
\[
 \begin{gathered}
 0\longrightarrow\mathcal C_S\xrightarrow{\ f_U\ }
             \Lambda_U\longrightarrow C\longrightarrow0,\\
 0\longrightarrow\operatorname{Tor}_1^S(C,B)
    \xrightarrow{\ \partial\ }\overline{\mathcal C}_S
    \xrightarrow{\ \overline f_U\ }\overline\Lambda_U
    \longrightarrow C\otimes_S B\longrightarrow0,\\
 \operatorname{Tor}_1^S(C,B)=\bigoplus_{r=p}^K B\epsilon_r,
       \qquad \partial(\epsilon_r)=\overline T^r,
       \qquad C\otimes_S B\simeq
                    \bigoplus_{r=p}^K B[\overline W_r].
 \end{gathered}
 \tag{TCL.19}\label{CFStage:eq:tcl-tor}
\]
Here \(\operatorname{Tor}_1\) can be computed directly, without
a flatness assertion about \(B\). The first line is a free
resolution of \(C\) of length one, whose differential in the
displayed bases is \(\Delta\). Tensoring this two-term complex
with \(B\) gives differential \(\overline\Delta\).
By definition its degree-one homology is
\(\operatorname{Tor}_1^S(C,B)\). The factorial calculation in
(TCL.16) gives its kernel with the displayed basis and gives its
cokernel with the displayed basis. The degree-one differential
was chosen as \(+f_U\), so the map of the kernel into the
degree-one free module is precisely
\(\epsilon_r\mapsto\overline T^r\), with no additional sign.
This proves the entire sequence and its maps explicitly.

For the top original vector the formula, including its full scalar, is
\[
 f_U(T^K)=\frac{K!}{((m-1)!)^k}\,a_0^k
                    \prod_{i=1}^k y_i^{m-1}.
 \tag{TCL.20}\label{CFStage:eq:tcl-terminal}
\]
Only \(\mathbf b=(m-1,\ldots,m-1)\) has total degree \(K\),
so (TCL.7) gives its multinomial coefficient. Every nonconstant
term of \(U\) annihilates that monomial; its constant coefficient
is exactly \(a_0^k\), proving the formula. Its scalar is nonzero
in the actual characteristic-zero domain, and its exact factorial
depth is \(d_K\), because \((m-1)!\) and \(a_0\) are units.
The displayed vector maps to zero under the specified characteristic
\(p\) coefficient map exactly when \(K\geq p\). This statement
concerns the already specified local single-primary source. It does
not extend that specialization to a larger coefficient ring containing
additional inverse CRT denominators. A ring in which \(p\) is a
unit cannot map unitally to characteristic \(p\), since applying
such a map to \(p p^{-1}=1\) would give \(0=1\).

\subsection{The exact original chain-support map for these depths}

Fix an integer \(L\geq\max(1,d_K)\) and put
\[
 \begin{gathered}
 R_{p,L}=\mathbb Z/p^L\mathbb Z,\qquad S_L=S/p^L S,\\
 \iota_L:R_{p,L}\longrightarrow S_L,
       \quad [z]\longmapsto[z],\\
 I_a=p^aR_{p,L},\quad H_a=p^aS_L\quad(0\leq a\leq L).
 \end{gathered}
 \tag{TCL.21}\label{CFStage:eq:tcl-chain-rings}
\]
The scalar map is well-defined and unital. It is injective.
Indeed, if an integer \(z\) not divisible by \(p^L\) has image
zero, write \(z=p^a u\) with \(0\leq a<L\) and \(u\) prime
to \(p\). Then \(p^a u=p^L s\) in \(S\). Cancellation and
inversion of \(u\) imply \(1=p^{L-a}s u^{-1}\), contradicting
the specified residue map. Thus precisely \(p^L\mathbb Z\)
maps to zero, proving injectivity. The same argument proves
\(p^a\notin p^{a+1}S_L\) for \(a<L\): a purported equality
lifts to \(p^a=p^{a+1}s+p^L t\), and cancellation gives
\(1=p s+p^{L-a}t\), again impossible in the residue field.
Consequently all \(H_a\) are distinct.

For completeness, every ideal of \(R_{p,L}\) is some \(I_a\).
For a nonzero ideal, choose a nonzero element with least exponent
of \(p\) among its integer representatives. It has the form
\(p^a u\) with \(u\) a unit modulo \(p^L\), by the integer
Euclidean algorithm. Multiplying by \(u^{-1}\) puts \(p^a\)
in the ideal, and all its other elements have exponent at least
\(a\) by choice, proving that the ideal is \(I_a\). The zero
ideal is \(I_L\). This argument is not applied to \(S_L\):
there we retain exactly the principal subchain \(\{H_a\}\).

Extension and contraction on this actual subchain are
\[
 \iota_L(I_a)S_L=H_a,\qquad \iota_L^{-1}(H_a)=I_a,
 \qquad(0\leq a\leq L).
 \tag{TCL.22}\label{CFStage:eq:tcl-extend-contract}
\]
Extension follows because the image of the principal generator is
\(p^a\). For contraction, an integer of exponent \(b<a\)
belonging to \(p^aS_L\) would give
\(p^b u=p^a s+p^L t\). Cancelling \(p^b\) contradicts the
nonzero residue of the unit \(u\). Thus only the integers divisible
by \(p^a\) contract, and they all do, proving the second formula.

The original ideal laws are retained exactly:
\[
 \begin{array}{lll}
 I_a+I_b=I_{\min(a,b)},& I_a\cap I_b=I_{\max(a,b)},&
              I_a I_b=I_{\min(L,a+b)},\\
 H_a+H_b=H_{\min(a,b)},& H_a\cap H_b=H_{\max(a,b)},&
              H_a H_b=H_{\min(L,a+b)}.
 \end{array}
 \tag{TCL.23}\label{CFStage:eq:tcl-ideal-laws}
\]
For sum and intersection the proof is inclusion: when \(a\leq b\),
the ideal indexed by \(b\) is contained in the one indexed by
\(a\). Their sum is the larger and their intersection the smaller.
For product, every sum of products of multiples of \(p^a\) and
\(p^b\) is a multiple of \(p^{a+b}\), and the generator
\(p^{a+b}\) itself is such a product. It becomes zero for
\(a+b\geq L\). These arguments apply in both displayed rings
and prove all six laws, including the boundary indices.

Let \(C_L=\{0<1<\cdots<L\}\). The original lattice map and its
actual coefficient extension are
\[
 \theta_p:C_L\longrightarrow\operatorname{Id}(R_{p,L}),
       \quad j\longmapsto I_{L-j},\qquad
 \theta_S:C_L\longrightarrow\{H_a:0\leq a\leq L\},
       \quad j\longmapsto H_{L-j}.
 \tag{TCL.24}\label{CFStage:eq:tcl-theta}
\]
They are bijections by the distinctness and classification just
proved. Each takes bottom to the zero ideal, top to the unit ideal,
maximum to ideal sum and minimum to ideal intersection, by
(TCL.23). Thus they are bounded-lattice isomorphisms onto their
stated codomains. Extension \(I_a\mapsto H_a\) commutes with
these two maps and preserves all three ideal operations of
(TCL.23). In the \(j\)-coordinates, ordinary ideal multiplication
is \(j\mathbin{\odot}t=\max(0,j+t-L)\), since its exponent is
\(\min(L,(L-j)+(L-t))\). It is not identified with lattice
meet: for \(L\geq2\), \(j=t=L-1\) gives respective values
\(L-2\) and \(L-1\).

For either coefficient ring \(B=R_{p,L}\) or \(B=S_L\), use
precisely the original split-support carrier
\[
 \begin{split}
 G_{C_L}(B)&=\{(0_B,j):0\leq j\leq L\}
                 \cup\{(b,L):b\in B\},\\
 (b,j)+(c,t)&=(b+c,\max(j,t)),\qquad
 (b,j)(c,t)=(bc,\min(j,t)),\\
 \tau_a&=(0_B,L-a)\ (0\leq a\leq L),\qquad
 e=\tau_0=(0_B,L),\qquad \tau_L=(0_B,0).
 \end{split}
 \tag{TCL.25}\label{CFStage:eq:tcl-carrier}
\]
The union is inside \(B\times C_L\), so the two occurrences
of \((0_B,L)\) give one element. A nonzero amplitude occurs only
at top support. A nonzero sum has a nonzero summand, so has top
support; a nonzero product has both amplitudes nonzero, so has
top support. This proves closure. Associativity and commutativity
hold in each coordinate. Distributivity follows from ring
distributivity and \(\min(j,\max(t,v))
=\max(\min(j,t),\min(j,v))\), verified by ordering \(t,v\).
The zero is \(\tau_L\), and the unit is \((1_B,L)\).
These facts prove the stated semiring structure directly.

The exact support-preserving coefficient morphism is
\[
 \begin{split}
 G_{C_L}(\iota_L):G_{C_L}(R_{p,L})&\longrightarrow G_{C_L}(S_L),
           \quad (b,j)\longmapsto(\iota_L(b),j),\\
 \chi_B:G_{C_L}(B)&\longrightarrow C_L,
           \quad(b,j)\longmapsto j,\\
 \theta_S\chi_{S_L}\,G_{C_L}(\iota_L)
    &=\operatorname{ext}_{\iota_L}\,\theta_p\chi_{R_{p,L}}.
 \end{split}
 \tag{TCL.26}\label{CFStage:eq:tcl-support-map}
\]
The first map is well-defined because nonzero target amplitude
has nonzero source amplitude and hence top support. Applying the
ring homomorphism to amplitudes and leaving the support coordinate
fixed proves preservation of both operations, zero and unit.
Its injectivity follows from injectivity of \(\iota_L\) and
the unchanged support coordinate. The support projection preserves
the two operations by their defining formulas. The last equality
follows on every entry \((b,j)\) from
\(I_{L-j}S_L=H_{L-j}\); it is an equality into the principal
ideal subchain with ideal sum and ideal intersection. In particular
each \(\tau_a\) maps to the distinct \(\tau_a\), then to
\(I_a\) or \(H_a\). The supported zero \(e\) remains distinct
from the absent zero \(\tau_L\).

Finally the connection of this support diagram to the actual cyclic
quotient is not an identification of a torsion module with a semiring.
The precise coordinate and annihilator maps are
\[
 \begin{split}
 C&\simeq\bigoplus_{r=0}^K S/(p^{d_r}),\qquad p^L C=0,\\
 \operatorname{Ann}_{S_L}(\nu_r)&=H_{d_r}
     =\theta_S(L-d_r)
     =\theta_S\chi_{S_L}(\tau_{d_r}),\\
 q_r(x)&=\lambda_r(U^{-1}x)\bmod p^{d_r}.
 \end{split}
 \tag{TCL.27}\label{CFStage:eq:tcl-quotient-support}
\]
The first and last lines are the proved maps of (TCL.15), with
\(p^L\) annihilating every summand because \(d_r\leq L\).
An element of \(S_L\) kills the unit class of the \(r\)-th
summand exactly when a representative in \(S\) lies in
\(p^{d_r}S\); this proves the annihilator formula and all its
equalities using (TCL.24)--(TCL.26). Thus the depths of the
original factorial inclusion have explicit labels in the original
chain-support carrier, with their arithmetic ideal product as well
as their lattice sum and intersection still recorded.
The entry \(\tau_{d_r}\) labels the annihilator; it does not
replace the vector \(\nu_r\). In particular, when \(d_r=L\)
the vector \(\nu_r\) is the nonzero unit class in \(S_L\), whereas
its annihilator is the zero ideal labelled by \(\tau_L\).

\subsection{The retained boundary-character labels and the full-stage obstruction}

The action on the lost source directions is also exact. For
\(K\geq p\), the kernel in (TCL.16), equivalently the connecting
image of the Tor group in (TCL.19), has the ordered basis
\(\overline T^p,\ldots,\overline T^K\), with
\[
 M_T\overline T^r=\overline T^{r+1}\ (p\leq r<K),\qquad
 M_T\overline T^K=0,\qquad
 \operatorname{ind}\bigl(M_T|_{\ker\overline f_U}\bigr)=K-p+1.
 \tag{TCL.28}\label{CFStage:eq:tcl-tor-index}
\]
These equalities follow in the free monomial source basis. In
particular the \((K-p)\)-th power takes \(\overline T^p\)
to the nonzero \(\overline T^K\), and the next power is zero
on every basis element. This index can exceed \(p\). It is an
action on the exact kernel in the source, not a subspace of the
ambient target. Thus it is compatible with
\(\overline J^p=0\) on \(\overline E\): the original map is
zero on every one of these kernel vectors.

Put \(N=m+1\). For a residue class \(c\in\mathbb Z/N\mathbb Z\)
define the original coefficient character masks in \(y\)-coordinates
and their full-unit transports by
\[
 \begin{split}
 Q_c(y^{\mathbf b})&=
       \mathbf 1_{\{|\mathbf b|+k\equiv c\ ({\rm mod}\ N)\}}
                       y^{\mathbf b},\qquad
 \widehat Q_c=UQ_cU^{-1}:E_S\longrightarrow E_S,\\
 D_c(T^r)&=\mathbf 1_{\{r+k\equiv c\ ({\rm mod}\ N)\}}T^r
                       \quad\hbox{on }\mathcal C_S.
 \end{split}
 \tag{TCL.29}\label{CFStage:eq:tcl-character-mask}
\]
Each diagonal entry defining \(Q_c,D_c\) is zero or one in
\(S\), so these maps exist over this exact coefficient ring
without adjoining roots of unity. On every indicated basis vector
exactly one of the masks is one. Hence each family consists of
pairwise orthogonal algebraic idempotents summing to the identity;
conjugation by the actual \(U\) proves the same assertion for
\(\widehat Q_c\). The word orthogonal here means that distinct
idempotents compose to zero; no inner product or orthogonal
projection on a metric space is introduced.

The full original boundary-character comparison is
\[
 \begin{split}
 Q_c V_r&=\mathbf 1_{\{r+k\equiv c\}}V_r,\qquad
 \widehat Q_c W_r=\mathbf 1_{\{r+k\equiv c\}}W_r,\\
 \widehat Q_c f_U&=f_U D_c,\qquad
 \widehat Q_c J=J\widehat Q_{c-1},\qquad
 D_c M_T=M_TD_{c-1}.
 \end{split}
 \tag{TCL.30}\label{CFStage:eq:tcl-character-diagram}
\]
Every summand of \(V_r\) has degree \(r\), proving the first
equality, and conjugation proves the second. Apply these formulas
to \(f_U(T^r)=r!W_r\) to prove the third equality on a source
basis. Multiplication by each \(y_i\) raises degree by one for
every surviving monomial and gives zero at the truncation boundary.
It follows on the complete monomial basis that
\(Q_cJ=JQ_{c-1}\). Conjugation and \(UJ=JU\) prove the fourth
identity; the identical degree argument on \(T^r\) proves the
last. These are covariant character shifts, not a claim that a
fixed character mask commutes with \(J\).

Over \(\mathbb C\), this is exactly the coefficient mask of
the boundary monodromy in TP.9--10: that period monodromy has
eigenvalue \(\exp(2\pi i(|\mathbf b|+k)/N)\) on its period
character indexed by \(\mathbf b\). In the unchanged original
\(s^{\mathbf b}\)-coordinates the coefficient mask is
\(P_k^{-1}Q_cP_k\), and the source mask is its conjugate by
\(P_k^{-1}[U]_{(y^{\mathbf b})}P_k\). The complex period matrix
is \(\Pi_k=e^{kc_0/u}W_kD_k(u)P_k\), with the original phase
constant here denoted \(c_0=-(-\rho)^N/N\), and \(D_k(u)\)
diagonal in these same \(\mathbf b\) indices. Consequently
\[
 \mathcal P_c\Pi_k=\Pi_k(P_k^{-1}Q_cP_k),\qquad
 \mathcal P_c\,\Pi_k[U]_{(s^{\mathbf b})}^{-1}
   =\Pi_k[U]_{(s^{\mathbf b})}^{-1}
                [\widehat Q_c]_{(s^{\mathbf b})},
 \quad \mathcal P_c=W_kQ_cW_k^{-1}.
 \tag{TCL.31}\label{CFStage:eq:tcl-period-mask}
\]
Indeed \(Q_c\) commutes with the diagonal \(D_k(u)\), which
proves the first equality by multiplication, and substitution of
the conjugation formula for \(\widehat Q_c\) proves the second.
The common scalar \(e^{kc_0/u}\), the original \(P_k\) and the
entire \(U\) are present in these identities.

The masks preserve both \(L_U\) and \(\Lambda_U\), as follows
from their diagonal action on \(r!W_r\) and \(W_r\). They
therefore act on the quotient \(C\), and the exact free resolution
in (TCL.19) commutes with \(D_c\) on its source and
\(\widehat Q_c\) on its target. Its base change and its kernel
inherit precisely those same masks. Thus the simultaneously retained
depth and character of every quotient coordinate and lost direction are
\[
 \begin{split}
 \nu_r=[W_r]&:\quad
       \operatorname{Ann}_{S_L}(\nu_r)=H_{d_r},\quad
       \widehat Q_c\nu_r=\mathbf 1_{\{r+k\equiv c\}}\nu_r,\\
 \epsilon_r\in\operatorname{Tor}_1^S(C,B)\quad(r\geq p)
     &:\quad \overline D_c\partial\epsilon_r
        =\mathbf 1_{\{r+k\equiv c\}}\overline T^r.
 \end{split}
 \tag{TCL.32}\label{CFStage:eq:tcl-two-labels}
\]
The first line combines the proved annihilator in (TCL.27) with
the induced quotient map. The second follows from the explicit
connecting map \(\partial\epsilon_r=\overline T^r\).
This proves both labels on the actual maps, including the zero
coordinates when \(d_r=0\). It does not assert an inertia-invariant
description of a different sheaf after reduction.

There is an exact obstruction to carrying this local specialization
through the full amplified quartet CRT stage. In the counterfactual
quartet of PAM.1--4 retain
\(\rho=\tfrac12+\delta+i\gamma\), with
\(0<\delta<\tfrac12\), \(\gamma>0\), and take \(k\geq p\).
Two actual ordered tensor centres in that construction are
\[
 \lambda_0=k\rho,\qquad
 \lambda_1=(k-p)\rho+p(1-\overline\rho),\qquad
 \lambda_0-\lambda_1=p(\rho+\overline\rho-1)=2p\delta\ne0.
 \tag{TCL.33}\label{CFStage:eq:tcl-crt-centres}
\]
The first chooses \(\rho\) in every factor; the second chooses
\(1-\overline\rho\) in exactly \(p\) factors and \(\rho\)
in the other \(k-p\) factors. Both choices are permitted in the
retained quartet, and the difference follows by subtraction.
In the grid coordinates of PAM.2 they have indices respectively
\((a,b)=(k,k)\) and \((k-p,k)\).

Let \(A_{\rm CRT}\subset\mathbb C\) be the actual coefficient
ring carrying those quartet coordinates and every coefficient of
the retained terminal CRT idempotent polynomial representative
\(e(Z)=\sum_{r=0}^{d}e_rZ^r\) at \(\lambda_0\).
By its specified CRT residues, \(e(\lambda_0)=1\) and
\(e(\lambda_1)=0\). Define the following element of this same ring:
\[
 \mathscr D=\sum_{r=1}^{d}e_r
       \sum_{j=0}^{r-1}\lambda_0^{r-1-j}\lambda_1^j
       \in A_{\rm CRT},\qquad
 1=p(\rho+\overline\rho-1)\mathscr D,
 \qquad
 p^{-1}=(\rho+\overline\rho-1)\mathscr D\in A_{\rm CRT}.
 \tag{TCL.34}\label{CFStage:eq:tcl-crt-obstruction}
\]
For every \(r\geq1\), multiplying the inner sum by
\(\lambda_0-\lambda_1\) cancels all intermediate monomials
and leaves \(\lambda_0^r-\lambda_1^r\). After multiplication
by \(e_r\) and summation the result is
\(e(\lambda_0)-e(\lambda_1)=1\). Substitution of (TCL.33)
proves the second equality and then the displayed inverse of
\(p\). This proof uses the original CRT polynomial itself and
does not depend on a selected formula for its coefficients or
on an assumed list of inverted pairwise differences.
Since \(p\) is a unit, no prime ideal of \(A_{\rm CRT}\)
lies over \((p)\), and there is no unital characteristic-\(p\)
specialization of that complete stage. A prime containing \(p\)
would contain \(1\), and any such specialization would send the
proved identity \(p p^{-1}=1\) to \(0=1\).

Thus for \(k\geq p\) the actual local prime used in (TCL.1)
does not extend to this full quartet stage with its retained
terminal CRT idempotent. This is a proved obstruction for that
specified stage, not a statement that the local and global
characteristic-zero objects are unrelated. Their local cyclic
coefficient map is still exactly (TCL.4), their unchanged monomial
translation is (TCL.13), and the characteristic-zero terminal
coefficient is (TCL.20). The local characteristic-\(p\) quotient
and its lost directions are the maps (TCL.16)--(TCL.19); a missing
extension of their coefficient map is not the image of a global
vector under a nonexistent specialization.

More generally let \(A_{{\rm CRT},k}\subset\mathbb C\) denote the
retained coefficient ring of the complete split corner at degree
\(k\). The argument (TCL.33)--(TCL.34), applied to every rational
prime \(q\leq k\) in place of \(p\), proves
\[
 q^{-1}\in A_{{\rm CRT},k}\quad(q\text{ prime},\ q\leq k).
 \tag{TCL.35}\label{CFStage:eq:tcl-all-small-primes}
\]
Indeed the second centre uses \(q\) reflected factors, which are
available for every such \(q\), and the same retained idempotent
separates it from the terminal centre. Its divided difference
gives the inverse without changing the idempotent.
For a directed coefficient ring containing all these actual
split stages one may take the explicit subring
\[
 A_{{\rm CRT},\infty}
   =\bigcup_{n\geq2}
       \mathbb Z[A_{{\rm CRT},2},\ldots,A_{{\rm CRT},n}]
       \ \subset\mathbb C,\qquad
 \mathbb Q\subset A_{{\rm CRT},\infty}.
 \tag{TCL.36}\label{CFStage:eq:tcl-split-limit}
\]
The finite generated-stage subrings in this union form an
increasing sequence, so the union is a subring containing every
stage with its actual coefficient maps. For any prime \(q\),
the stage \(k=q\) supplies \(q^{-1}\) by (TCL.35).
Factoring a nonzero integer into primes then puts its reciprocal
in the union, proving the asserted copy of \(\mathbb Q\).
Every residue field of this union therefore has characteristic
zero: a characteristic-\(q\) ring homomorphism would send the
unit identity \(q q^{-1}=1\) to \(0=1\).
Equivalently the scalar morphism
\(\operatorname{Spec}A_{{\rm CRT},\infty}
 \to\operatorname{Spec}\mathbb Z\)
has no point over any \((q)\). This follows directly by contraction
of primes and the proved invertibility of every \(q\).
These conclusions concern the full split corner diagram with
its retained CRT idempotents. They do not remove the finite
single-primary models or the local maps (TCL.1)--(TCL.32).

There are also exact surviving character-zero vectors in the
reduced local cyclic image. Take \(m\geq2\) and \(k=p\ell\) with
\(\ell\geq1\), so \(K=k(m-1)\geq p\). Its final nonzero cyclic
image basis vector and its unchanged characteristic-zero precursor
satisfy
\[
 \begin{split}
 \overline J\,\overline W_{p-1}&=0,\qquad
       \overline W_{p-1}\ne0,\qquad J W_{p-1}=pW_p\ne0,\\
 \overline{\widehat Q}_0\,\overline W_{p-1}=\overline W_{p-1}
       &\quad\Longleftrightarrow\quad
       p(\ell+1)\equiv1\pmod{N},\qquad N=m+1 .
 \end{split}
 \tag{TCL.37}\label{CFStage:eq:tcl-surviving-character}
\]
The first line follows from (TCL.12), the unit pivot after
reduction, and the nonzero scalar \(p\) on the free
characteristic-zero containing lattice. The vector belongs to
the original reduced cyclic image because
\(\overline f_U(\overline T^{p-1})
 =\overline{(p-1)!}\,\overline W_{p-1}\), whose scalar is a
unit. Formula (TCL.30) gives the second line as
\(p-1+k\equiv0\pmod N\), and substitution of \(k=p\ell\)
gives the stated congruence. In the retained coefficient datum
\(p>N\), so \(p\) is invertible modulo \(N\). Choosing the
single residue class
\(\ell\equiv p^{-1}-1\pmod N\) and then its infinitely many
positive representatives proves that the displayed character-zero
case occurs for infinitely many tensor orders \(k=p\ell\).
These vectors are killed by \(J\) only after base change, as
their exact characteristic-zero image \(pW_p\) shows. This is
a calculation in the local coefficient module with its retained
character masks, not an identification with a finite-field
sheaf vector or a recovery of the original terminal jet.

Every tensor coefficient, the original monomial translations, the
full unit and inverse, and the specified coefficient map occur in
the displayed morphisms. The characteristic-\(p\) kernel in
(TCL.16) is a statement about these coefficient modules. It neither
sets an integer scalar to zero in \(\overline{\mathbb Q}_\ell\)
nor identifies this reduction with the inertia or Frobenius action
on the finite-field sheaf. No verdict about the location of an
actual zero of \(\xi\) follows from this lattice calculation alone.

% END TCL
% END COMPLETE TENSOR COMPANIONS
